% concatenated from main.tex, MP_picture.tex
%%%%%%%%%% HEADER STUFF %%%%%%%%%%
\documentclass{article}
\usepackage[hmargin=2.5cm,vmargin=3cm,bindingoffset=0.5cm]{geometry}


\usepackage{fancyhdr,graphics,graphicx,xy}
\usepackage{ifthen}
\usepackage{amssymb,amsmath,mathtools,enumerate,tikz-cd}
\usepackage{amsthm}
\usepackage{graphics}
\usepackage{dsfont}
\usepackage{algpseudocode}
\usepackage{adjustbox}
\usepackage{bbm}
\usepackage{tikz}
\usepackage[affil-it]{authblk}
\usepackage{circuitikz}
\usepackage[permil]{overpic}


\usetikzlibrary{arrows}
\usetikzlibrary{snakes}
\graphicspath{/Users/huubdejong/Downloads/}
%\pagestyle{fancy}
\newcommand{\soltop}{  
\fancyhead[C]{Thesis notes}
}
%%%%%%%%% END HEADER STUFF %%%%%%%%%





%%%%%%%%%%%% SHORTCUTS FOR IMPORTANT SYMBOLS %%%%%%%%%%%%

\def\N{\mathbb{N}}
\def\Q{\mathbb{Q}}
\def\C{\mathbb{C}}
\def\R{\mathbb{R}}
\def\P{\mathbb{P}}
\def\A{\mathbb{A}}
\def\Z{\mathbb{Z}}
\def\E{\mathbb{E}}
\def\T{\mathbb{T}}

\def\F{\mathcal{F}}
\def\M{\mathcal{M}}
\def\mH{\mathcal{H}}
\def\mM{\M}
\def\mS{\mathcal{S}}
\def\mJ{\mathcal{J}}

\def\sumk{\sum_{k=1}^\infty}
\def\sumn{\sum_{n=1}^\infty}
\def\sumi{\sum_{i=1}^\infty}
\def\unionk{\bigcup_{k=1}^\infty}
\def\unionn{\bigcup_{n=1}^\infty}
\def\unioni{\bigcup_{i=1}^\infty}
\def\interi{\bigcap_{i=1}^\infty}
\def\intern{\bigcap_{n=1}^\infty}
\def\interk{\bigcap_{k=1}^\infty}
\def\limn{\lim_{n\to\infty}}
\def\limk{\lim_{k\to\infty}}
\def\limj{\lim_{j\to\infty}}
\def\supp{\operatorname{supp}}
\def\diam{\operatorname{diam}}
\def\lcm{\operatorname{lcm}}

\newcommand{\brac}[1]{\left({#1}\right)}
\newcommand{\cbrac}[1]{\left\{{#1}\right\}}
\newcommand{\abs}[1]{\left\vert{#1}\right\vert}
\newcommand{\Set}[1]{\left\{ #1 \right\}}
\newcommand{\set}[1]{\{ #1 \}}
\newcommand{\mtr}[1]{\begin{pmatrix} #1 \end{pmatrix}}

\newcommand{\ind}[1]{\textbf{1}_{\cbrac{{#1}}}}

\newcommand{\norm}[1]{\left\|{#1}\right\|}
\newcommand{\gnorm}{\left\|\cdot\right\|}

% \DeclareMathOperator{\lcm}{lcm}
\DeclareMathOperator{\proj}{proj}

\title{Smoothness of Markov Partitions for Expanding Toral Endomorphisms}
\author{Chayce Hughes%
  \thanks{Electronic address: \texttt{chahug21@student.ubc.ca}}\;}
\author{Huub de Jong%
  \thanks{Electronic address: \texttt{huub.dejong@math.ubc.ca}}}

\affil{The University of British Columbia, Vancouver}
\date{\today}

%%%%%%%%%% THEOREM COMMANDS %%%%%%%%%%
\newtheorem{theorem}{Theorem}[section]
\newtheorem{corollary}[theorem]{Corollary}
\newtheorem{definition}[theorem]{Definition}
\newtheorem{lemma}[theorem]{Lemma}
\newtheorem{example}[theorem]{Example}
\newtheorem{proposition}[theorem]{Proposition}
\newtheorem{remark}[theorem]{Remark}
\newtheorem{conjecture}[theorem]{Conjecture}






\begin{document}

\maketitle

\begin{abstract}
    We show that an expanding toral endomorphism in dimension 2 admits a smooth (in fact linear) Markov partition if and only if some power of the corresponding integer matrix is diagonalizable with integer eigenvalues. We exhibit examples of qualitatively different smoothness behavior, and highlight the existence of a ``hybrid'' type of smoothness in dimension 2. For dimension $d$, we show that expanding toral endomorphisms satisfying the eigenvalue condition above admit a linear Markov partition. Finally, we provide an estimate on the Hausdorff dimension of the boundary of a Markov partition using techniques from symbolic dynamics.
\end{abstract}

\section{Introduction} \label{Section:Intro}

Markov partitions were developed by Adler and Weiss for hyperbolic toral automorphisms \cite{AW}, and systematically studied by Sinai in the context of Anosov diffeomorphisms on manifolds \cite{Si}. The utility of Markov partitions comes from the fact that they give rise to a symbolic description of the system which retains large amounts of interesting information. A great deal of advancment in the theory on Markov partitions is due to Bowen, who used these partitions as a tool in better understanding the more general class of Axiom A diffeomorphisms \cite{Bo1}. He showed that such diffeomorphisms admit a Markov partition and used them to show several dynamical properties, such as the 0-dimensionality of minimal sets \cite{Bo3}. 

A particularly beautiful 3-page paper of Bowen \cite{Bo2} shows that no Markov partition for a hyperbolic toral automorphism in dimension 3 can have a piecewise smooth boundary. This result was later extended by Cawley \cite{C} to completely classify when a hyperbolic toral automorphism admits a Markov partition with piecewise smooth boundary, and to show that a piecewise smooth boundary must in fact be linear. 

\begin{theorem}[Cawley, 1991]\label{thm:Cawley}
    A hyperbolic toral automorphism $f:\T^d\to\T^d$ induced by a matrix $A$ admits a smooth Markov partition if and only if some power of $A$ is similar over $\Q$ to a block diagonal matrix $$\begin{pmatrix}
        L_1 & & & \\
        & L_2 & & \\
        & & \ddots & \\
        & & & L_{d/2}
    \end{pmatrix}$$
    where each $L_i$ is a $2\times 2$ hyperbolic integer matrix with $\det L_i = \pm1$. 

    Moreover, any smooth Markov partition for a hyperbolic toral automorphism must be linear.
\end{theorem}

In general, Markov partitions are developed for invertible systems, but the theory applies in the context of non-invertible, expanding dynamical systems as well. The existence of Markov partitions in such systems is shown in e.g. Chapter 4 of \cite{URM} or Chapter 7 of \cite{Ru}, and a side-by-side development of the invertible and non-invertible settings can be found in \cite{AKS}. Work of Bedford \cite{Be} uses explicitly fractal techniques due to Dekking \cite{D} to construct Markov partitions for expanding toral endomorphisms. A folkloric result says that these Markov partitions ``almost always'' have non-smooth boundary, but we were not able to locate a precise statement in the literature.

This work attempts to address this apparent gap in the literature, and proves some smoothness results that seem to parallel Theorem \ref{thm:Cawley}, although we obtain them by elementary tools, avoiding the heavy topological machinery used in \cite{C}. More precisely, we prove the following; all the relevant definitions, in particular those describing different types of smoothness, can be found in Section 2. 



\begin{theorem}\label{thm:power diag with integer eigenvals implies smooth MP}
    Suppose $A$ is a $d\times d$ expanding integer matrix such that for some $k\in\N$, $A^k$ is diagonalizable with integer eigenvalues. Then the induced expanding toral endomorphism $f$ admits a linear Markov partition.
\end{theorem}

\begin{theorem}\label{thm:nonsmoothness in dimension 2}
    Let $A$ be a $2\times 2$ expanding integer matrix, $f$ the induced expanding toral endomorphism, and $\M$ a Markov partition for $f$. Then the following are true: \begin{enumerate}
        \item If $A$ has real, irrational eigenvalues of different modulus, then $\M$ is essentially nowhere smooth.
        \item If $A$ has complex eigenvalues with arguments that are not a rational multiple of $\pi$, then $\M$ is nowhere differentiable.
        \item If $A$ is not diagonalizable, then $\M$ is not smooth, but may be hybrid.
    \end{enumerate}
\end{theorem}

Theorem \ref{thm:Jordans have hybrid MP} will show that in fact every non-diagonalizable expanding $2\times 2$ matrix admits a hybrid Markov partition. Combining the two theorems above and the Rational Root Theorem gives us the following.
\begin{corollary}
    Let $A$ be an expanding $2\times 2$ integer matrix. Then the induced expanding toral endomorphism $f$ admits a smooth, in fact linear, Markov partition if and only if for some $k\in \N$, $A^k$ is diagonalizable with integer eigenvalues.
\end{corollary}

There is an obvious conjectured generalization of this result, mirroring Theorem \ref{thm:Cawley}, which we leave for future work.

\begin{conjecture}
    Let $A$ be an expanding $n\times n$ integer matrix. Then the induced toral endomorphism $f$ admits a smooth, in fact linear, Markov partition if and only if for some $k\in \N$, $A^k$ is diagonalizable with integer eigenvalues.
\end{conjecture}

The structure of this paper is as follows. Section \ref{Section:Prelims} contains preliminaries and a brief introduction to symbolic dynamics. Section \ref{Section:Constructions} is dedicated to proving Theorem \ref{thm:power diag with integer eigenvals implies smooth MP}, and gives some techniques of transforming a Markov partition for one system into a Markov partition for another. We also show that the power in the theorem statement is bounded, and expand on a more precise classification of all possible powers. In Section \ref{Section:NonSmooth}, we prove Theorem \ref{thm:nonsmoothness in dimension 2} via a case-by-case analysis. Section \ref{Section:FractalDim} is dedicated to estimating the Hausdorff dimension of the boundary of Markov partitions for arbitrary expanding systems. 

\section*{Acknowledgements}

The authors thank Mariusz Urbánski for providing background, advice, and encouragement. They also give thanks to Brian Marcus, who helped us greatly simplify our proofs, and Chengyu Wu and Alexia Yavicoli for stimulating discussion and detailed feedback. 

\section{Preliminaries} \label{Section:Prelims}


A \textit{dynamical system} is a compact metric space $X$ with a continuous self-map $f:X\to X$. The map $f:X\to X$ is \textit{positively expansive} if there exists some $c>0$, called an \textit{expansive constant},  such that for all $x,y\in X$, $x\neq y$ implies there exists an $n\in \N$ such that $d(f^nx,f^ny)>c$. We are interested in \textit{expanding toral endomorphisms} on the $d$-dimensional torus $\T^d:= \R^d/\Z^d$, which are maps induced from a $d\times d$ integer matrix $A$ for which all eigenvalues have modulus strictly bigger than 1. The induced expanding toral endomorphism is $$f_A:\T^d\to\T^d$$ defined as $$f_A(x+\Z^d)=Ax+\Z^d.$$ It is not hard to show that the condition on the modulus of the eigenvalues implies that $f_A$ is positively expansive. We will call $A$ an \textit{expanding} integer matrix and write $f$ for $f_A$ when no ambiguity exists. Write $p:\R^d\to\T^d$ for the standard projection map, and consider the standard metric on $\T^d$, $$d(x,y):=\min_{\substack{u,v\in \R^d \\ p(u)=x \\p(v)=y}} d_E(u,v)$$ where $d_E(u,v)$ is the Euclidian metric on $\R^d$.

Since $f:\T^d\to \T^d$ is a positively expansive map from a manifold onto itself, a theorem of Coven and Reddy \cite{CR} gives a metric $\tilde{d}$ compatible with the standard metric for which there exist constants $\varepsilon>0,\lambda >1$, so that \begin{equation*}
    \tilde d(x,y)<\varepsilon \implies \tilde d(fx,fy)>\lambda \tilde d(x,y)
\end{equation*} 
and \begin{equation*}
    \tilde d(fx,z)<\varepsilon\lambda \implies B(x,\varepsilon)\cap f^{-1}(z) \text{ is a singleton.}
\end{equation*} That is, $f$ is \textit{expanding} in their terminology with respect to $\tilde{d}$, which need not be equal to $d$. 

A set $R\subset X$ is \textit{proper} if $R=\overline{\operatorname{int}(R)}\neq \emptyset$. Following \cite{AKS}, a \textit{Markov partition} for an expansive dynamical system $(X,f)$ is a finite cover $\M=\{R_i\}_{i=1}^n$ of $X$ by proper sets with disjoint interiors such that two things hold: \begin{enumerate}
    \item $\M$ is \textit{metrically generating}, meaning for all $(x_i)\in \{1,...,n\}^\N$, $$\bigcap_{k=0}^\infty f^{-k}(\operatorname{int}(R_{x_n}))$$ consists of at most one point. 
    \item $\M$ satisfies a \textit{Markov property} (MP), meaning that for all $R,R'\in \M$, if $$f(\operatorname{int}(R))\cap \operatorname{int} (R')\neq \emptyset,$$ then each $y\in \operatorname{int} (R')$ has a unique preimage in $R$. 
\end{enumerate}

For constructive arguments, it is convenient to use a slightly stronger version of the Markov property above: 

\begin{enumerate}[(MP+)]
    \item For all $R\in \M$, $f|_R$ is injective, and there exist $R_1,...,R_n\in \M$ such that $f(R)=\bigcup_{i=0}^nR_i$.
\end{enumerate} The property (MP+) directly implies (MP), but it also captures more explicitly the intuitive idea behind Markov partitions for non-invertible systems. 

Existence of Markov partitions is commonly shown for systems $(X,f)$ which are expanding with respect to a given metric, which is why we point out the existence of $\tilde d$ and the expanding property. However, the Markov partition itself is a purely topological concept and makes no reference to the metric. Except for Section 5, we will always take any distance under consideration with respect to the standard metric $d$ on $\T^d$ or the Euclidian metric on $\R^d$. When no ambiguity exists, either metric will be denoted $d$. For details on Markov partitions for expanding systems in generality, we refer to Chapter 4 of \cite{URM} or Chapter 7 of \cite{Ru}. 



We denote the \textit{boundary} of a Markov partition $$\partial\M:=\bigcup_{R\in \M} \partial R.$$ It follows immediately from the definition that the boundary of a Markov partition is an invariant set.

\begin{lemma}\label{lem:boundaries are invariant}
    Let $f:X\to X$ be a positively expansive system on a compact metric space, and $\M$ a Markov partition for $f$. Then $$f(\partial\M)\subset \partial\M.$$
\end{lemma}

\begin{proof}
    Let $x\in \partial R$ and suppose $f(x)\in \operatorname{int}(R')$ for some $R,R'\in \M$. Moreover, let $c>0$ be an expansive constant for $f$. For all $n\in \N$ sufficiently large, take $\delta_n\in(0,c/2)$ so that $$f(B(x,\delta_n))\subset B\left(f(x),\frac{1}{n}\right) \subset \operatorname{int}(R').$$
    Since $x\in \partial R$, there exists some $\tilde{x}_n\in B(x,\delta_n)\setminus R$. By the Markov property, $f(\tilde{x}_n)\in \operatorname{int}(R')$ has a unique preimage $z_n$ in $R$, but $z_n$ cannot be in $B(x,c/2)$, for otherwise $z_n=\tilde{x}_n$ by positive expansiveness. Take now a convergent subsequence $z_{n_k}$ with $\bar{z}=\limk z_{n_k}\in R$, since $R$ is closed. Then $f(\bar{z})=\limk f(z_{n_k}) = f(x)$, and $\bar{z}\neq x$ since $d(x,\bar{z})\geq c/2$. But this contradicts the Markov property, so we are done. 
\end{proof}

We will classify qualitatively different types of smoothness for boundaries of Markov partitions.

\begin{definition}\label{def:smoothness MP}
    Let $f:\T^d\to \T^d$ be an expanding toral endomorphism, and $\M$ a Markov partition for $f$. We say $\M$ is \begin{itemize}
    \item \textit{linear} if $\partial\M$ is piecewise linear,
    \item \textit{smooth} if $\partial\M$ is piecewise $C^1$, 
    \item \textit{hybrid} if $\partial\M$ is not piecewise $C^1$, but contains a $C^1$ arc,
    \item \textit{essentially nowhere smooth} if $\partial\M$ contains no $C^1$ arcs, 
    \item \textit{nowhere differentiable} if for every continuous injection $\gamma:[0,1]\to\partial\M$ and every $t\in[0,1]$, either $\gamma$ is not differentiable at $t$ or $\gamma'(t)=0$. 
\end{itemize}
\end{definition}


We will repeatedly use the following property of the boundary of Markov partitions, which is immediate from Lemma \ref{lem:boundaries are invariant}.
\begin{lemma}\label{lem:boundary not dense}
    If $\M$ is a Markov partition for $f$, then $\bigcup_{k=0}^\infty f^k(\partial\M) = \partial \M$ is nowhere dense.
\end{lemma}

\begin{proof}
    Suppose that for some $R\in\M$, $\partial R$ is dense in some open set $U\subset \T^d$. Then one would have that $\overline{R}\cap\overline{R^c}$ contains an open set, but this would mean there are two rectangles whose interiors intersect, a contradiction. 
\end{proof}

The reason Markov partitions are of great interest in dynamical systems is that they allow for the construction of a \textit{symbolic coding} of a system. We briefly introduce this coding here, focusing on the aspects that are relevant for our results; for more details we refer to \cite{AKS} or Chapter 6 of \cite{LM}. 

A given Markov partition $\M$ for a dynamical system $(X,f)$ defines an $n\times n$ matrix $T$ where $$T(i,j)=\begin{cases}
    1,& \text{if } f(\operatorname{int}(R_i))\cap \operatorname{int}(R_j) \neq \emptyset \\
    0,& \text{otherwise}
\end{cases}.$$ Such a matrix $T$ in turn defines a \textit{shift of finite type} (SFT) 

$$\Sigma_T:=\{(x_0x_1x_2...)\in \{1,...,n\}^\N:\forall i\in\N, T(x_i,x_{i+1})=1\},$$ which we equip with the shift map $\sigma$ defined by $$\sigma(x)_i=x_{i+1}.$$ We give $\{1,...,n\}$ the discrete topology and $\{1,...,n\}^\N$ the corresponding product topology. Then $\Sigma_T$ is compact in the subspace topology, $\sigma$-invariant, and metrizable. 

Observation 1.4 of \cite{AKS} shows that a Markov partition yields a continuous and surjective map $\pi:\Sigma_T\to X$ such that the diagram
$$\begin{tikzcd}
\Sigma_T \arrow[r, "\sigma"] \arrow[d, "\pi"] & \Sigma_T \arrow[d, "\pi"] \\
X \arrow[r, "f"]                              & X                        
\end{tikzcd}$$ commutes. The map $\pi$ described there, although not presented in this form, is defined as $$\pi(x):=\bigcap_{N=0}^\infty \overline{\bigcap_{i=0}^Nf^{-i}\operatorname{int}(R_{x_i})},$$ where we identify the intersection with the (necessarily unique) point it contains. The map $\pi$ is called a \textit{coding} map.

Using the coding, one can \textit{refine} Markov partitions to make them arbitrarily small (in any compatible metric). Continuity of $\pi$ and compactness of $\Sigma_T$ shows that taking $N$ sufficiently large, the diameters of the sets $$\overline{\bigcap_{i=0}^Nf^{-i}\operatorname{int}(R_{x_i})}$$ can uniformly be made arbitrarily small. Taking these sets to be small enough, it can be shown that they satisfy (MP+), and hence form a new Markov partition. Crucially for our purposes, this type of \textit{recoding} does not change the geometry of the boundaries in an impactful way.



A particular construction of Markov partitions for expanding endomorphisms of $\T^2$ using fractals can be found in \cite{Be}, which we will refer to later. 





\section{Constructing smooth Markov partitions} \label{Section:Constructions}

We first construct some Markov partitions, and discuss how a Markov partition for one expanding matrix can be transformed into a Markov partition for another. In order to obtain a Markov partition for an expanding toral endomorphism $f_A$, it suffices to construct a certain $\Z^d$-invariant \textit{tiling} $\mathcal{C}$ of $\R^d$. That is, a cover of $\R^d$ by proper compact sets with disjoint interiors, with the following properties: \begin{enumerate}
    \item \textit{(Translation invariance)} For any $v\in \Z^d$, $\mathcal{C}+v=\mathcal{C}$.
    \item \textit{(Finiteness)} $\#\{C\in \mathcal{C}:C\cap[0,1]^d\neq \emptyset\}<\infty.$
    \item \textit{(Smallness)} For all $C\in\mathcal{C}$, $\diam(C)<1$ and $\operatorname{diam}(AC)<1.$
    \item \textit{(Markov property)} For all $C\in\mathcal{C}$, there exist $C_1,...,C_n\in \mathcal{C}$ satisfying $$AC=\bigcup_{1\leq i\leq n} C_i.$$ 
\end{enumerate}

The first two conditions imply that the latter two need only be checked for all of the finitely many $C\in \mathcal{C}$ that intersect the unit cube. We will call a tiling with the properties above a \textit{Markov tiling} for $A$, and refer to the elements of $\mathcal{C}$ as \textit{tiles}. Up to a refinement, Markov partitions and Markov tilings stand in one-to-one correspondence with each other. 

\begin{proposition}\label{prop:MPs give Markov tilings}
    Let $A$ be a $d\times d$ expanding integer matrix. Then the projection of any Markov tiling for $A$ is a Markov partition for $f_A$. 

    Conversely, any Markov partition for $f_A$ can, after a refinement, be lifted to a Markov tiling of $\R^d$.
\end{proposition}

\begin{proof}
    First, suppose $\mathcal{C}$ is a Markov tiling for $A$. Then, by translation invariance and the smallness condition, $\operatorname{int}(pC)\cap \operatorname{int}(pC') \neq \emptyset$ if and only if $C=C'+v$ for some $v\in \Z^d$. Finiteness ensures that $p\mathcal{C}$ consists of finitely many distinct elements, and the smallness property implies in particular that $f|_{p(C)}$ is injective for any $C\in \mathcal{C}$. Finally, the Markov property for tilings means that $p\mathcal{C}$ satisfies (MP+), hence is a Markov partition. It remains to show that $p\mathcal{C}$ is metrically generating. Enumerate the elements of $p\mathcal{C}$ so that $$p\mathcal{C}=\{R_i\}_{i=1}^n,$$ and let $(x_i)_{i\in \N}\in \{1,...,n\}^{\N}$. Suppose $$x,y\in \bigcap_{i=0}^\infty f^{-i}(\operatorname{int}R_{x_i}),$$ and let $\bar{x},\bar{y}\in \R^d$ be two arbitrary lifts of $x,y$, and $C_{x_i}$ an arbitrary tile in the preimage of $R_{x_i}$. Then for all $i$, there exists some $v_i,w_i\in \Z^d$ such that $$\bar{x}+v_i,\bar{y}+w_i\in A^{-i}C_{x_i}.$$ Since $A$ is expanding, all eigenvalues of $A^{-1}$ have modulus less than 1. Since the tiles of $\mathcal{C}$ are uniformly bounded in diameter, we get $$\operatorname{diam}(A^{-i}C_{x_i})\xrightarrow[i\to \infty]{} 0.$$ This means $\lim_{i\to\infty} d(\bar{x}+v_i,\bar{y}+w_i)= 0,$ meaning that $$d(p\bar{x},p\bar{y})=d(x,y)=0,$$ so $x=y$. We conclude that $p\mathcal{C}$ is a Markov partition satisfying (MP+).

    Now suppose $\mathcal{M}$ is a Markov partition for $f_A$. Note that for $x,y\in\T^d$ sufficiently close, $$ d(f_Ax,f_Ay) \leq \|A\| d(x,y),$$ where $\|A\|$ denotes the operator norm of $A$ induced by the Euclidian norm on $\R^d$. Thus, we can refine $\M$ such that for all $R\in\mathcal{M}$, $$\diam(f_AR)<\frac{1}{2}\quad \text{and} \quad \diam(R)<\frac{1}{2}.$$

    For each $R_i$, pick an arbitrary point $x_i\in R_i$ and $\bar{x}_i\in \R^d$ some lift. Then for each $v\in \Z^d$ and each $y\in R_i$, there exists a unique $\bar{y}\in p^{-1}(y)\cap B(\bar{x}_i+v,1/2).$ We define our tiling to be $$\mathcal{C}:=\{p^{-1}R_i\cap B(\bar{x}_i+v,1/2):R_i\in\M,v\in\Z^d\}.$$
    That is, $\mathcal{C}$ is the collection of unique copies of each $R_i$ within 1/2-balls centered at $\Z^d$ translations of $\bar{x}_i$.
    Translation invariance of $\mathcal{C}$ is immediate, and the finiteness property follows from the fact that $\mathcal{M}$ is finite. Since $\diam(fR_i)<1/2$, we also have $\diam(AC)<1/2<1$ for all $C\in \mathcal{C}$. It remains to show that $AC=\bigcup_jC_j$ for some finite set $\{C_j\}\subset\mathcal{C}$. 

    First note that $$pAC=fpC=\bigcup_{j=1}^kR_j$$ for some $R_j\in \M$. Since $\diam(AC)<1/2$, for each $R_j$ there is a unique $v_j\in\Z^d$ such that $\bar{x}_j+v_j\in AC$. Thus, $AC$ consists of exactly those unique copies of $R_j$ containing $\bar{x}_j+v_j$, hence $AC$ is a finite union of tiles $C_j\in \mathcal{C}$. 
\end{proof}

To prove Theorem \ref{thm:power diag with integer eigenvals implies smooth MP}, we will use a short sequence of lemmas.

\begin{lemma}\label{lem:MP for diagonal}
    Let $D$ be a $d\times d$ expanding diagonal integer matrix. Then $D$ admits a linear Markov partition.
\end{lemma}

\begin{proof}
    Let $K\geq |\det D|$ such that $$\operatorname{diam}(D[0,1/K]^d)<1.$$ Then $$\mathcal{C}=\{[i_1/K,(i_1+1)/K]\times ... \times [i_d/K,(i_d+1)/K]:i_1,..,i_d\in\Z\}$$
    is easily shown to be a Markov tiling for $D$. 
\end{proof}

\begin{lemma}\label{lem:MPforsimilarmatrices}
    Let $A,B$ be $d\times d$ expanding integer matrices that are similar over $\Q$. Then for any Markov tiling $\mathcal{C}$ for $A$ there exists a linear transformation turning $\mathcal{C}$ into a Markov tiling for $B$.
\end{lemma}

\begin{proof}
    Since $A,B$ are similar over $\Q$, there exists an integer matrix $T$ invertible over $\Q$ such that $B=TAT^{-1}$, where $T^{-1}$ need not be an integer matrix. Let $K$ be a large enough multiple of $|\det T|$ such that $\frac{1}{K}T\mathcal{C}$ satisfies the smallness property, i.e. $$\operatorname{diam}\left(\frac{1}{K}TC\right)<1,\text{ and }\operatorname{diam}\left(B\left(\frac{1}{K} TC\right)\right)<1$$ for all of the finitely many $C\in \mathcal{C}$ for which $\frac{1}{K}TC$ intersects the unit cube. 
    
    For any $C\in \mathcal{C}$, we have $$B\left(\frac{1}{K} T\mathcal{C}\right)=\frac{1}{K} TAT^{-1}TC=\frac{1}{K}T\bigcup_{i}C_i=\bigcup_i\frac{1}{K}TC_i,$$ so $\frac{1}{K}T\mathcal{C}$ satisfies the Markov property for tilings.
    Moreover, for any $v\in \Z^d$, $$\frac{1}{K}T\mathcal{C}+v=\frac{1}{K}T(\mathcal{C}+KT^{-1}v).$$
    Since $|\det T|$ divides $K$, $KT^{-1}v\in \Z^d$ by Cramer's rule. $\mathcal{C}$ is translation invariant, so $$\frac{1}{K}T\mathcal{C}+v=\frac{1}{K}T\mathcal{C}.$$

    We conclude that $\frac{1}{K}T\mathcal{C}$ is a Markov tiling for $B$, as desired. 
\end{proof}


\begin{lemma}\label{lem:MPforpowers}
    Let $\mathcal{C}$ be a Markov tiling for $A^k$. Then up to a linear rescaling, the collection $\mathcal{E}$ of nonempty sets $$\overline{\operatorname{int}C_0\cap \operatorname{int}AC_1\cap...\cap\operatorname{int}A^{k-1}C_{k-1}}$$ for $C_i\in\mathcal{C}$ is a Markov tiling for $A$. 
\end{lemma}

\begin{proof}
    We first replace $\mathcal{C}$ with $\frac{1}{K|\det A|^k}\mathcal{C}$ where $K$ is large enough so that $\mathcal{E}$ satisfies the smallness condition. This way, we guarantee that $A^j\mathcal{C}$ is translation invariant for $j\in\{0,1,...,k-1\}$. It is immediate that $\mathcal{E}$ satisfies the finiteness property and covers $\R^d$ by proper sets with disjoint interiors. 

    Since translation by $v\in\Z^d$ is a homemorphism on $\R^d$, \begin{align*}
        &\overline{\operatorname{int}C_0\cap \operatorname{int}AC_1\cap...\cap\operatorname{int}A^{k-1}C_{k-1}} +v \\
        &= \overline{\operatorname{int}(C_0+v)\cap \operatorname{int}(AC_1+v)\cap...\cap\operatorname{int}(A^{k-1}C_{k-1}+v)} \\
        &= \overline{\operatorname{int}C_0'\cap \operatorname{int}AC_1'\cap...\cap\operatorname{int}A^{k-1}C_{k-1}'}
    \end{align*}
    where $C_i' \in \mathcal{C}$ for $i\in\{0,1,...,k-1\}$, so $\mathcal{E}$ is translation invariant.

    For the Markov property, we note that since $A$ is a homeomorphism on $\R^d$, we get

    \begin{align*}
        A\left(\overline{\operatorname{int}C_0\cap \operatorname{int}AC_1\cap...\cap\operatorname{int}A^{k-1}C_{k-1}}\right)&=\overline{\operatorname{int} AC_0\cap \operatorname{int} A^2C_1\cap... \cap \operatorname{int}A^kC_{k-1}} \\
        &= \overline{\operatorname{int} AC_0\cap \operatorname{int} A^2C_1\cap... \cap \operatorname{int}\left(\bigcup_{1\leq j\leq n} C_j\right)} \\
        &= \bigcup_{1\leq j\leq n} \overline{\operatorname{int}C_j\cap\operatorname{int}AC_0\cap...\cap \operatorname{int}A^{k-1}C_{k-2}}
    \end{align*}
    
\end{proof}



Observe that the operations described in Lemmas \ref{lem:MPforsimilarmatrices} and \ref{lem:MPforpowers} preserve linearity, smoothness (or lack thereof), and the Hausdorff dimension of Markov partitions. Our first result is now immediate.

\begin{proof}[Proof of Theorem \ref{thm:power diag with integer eigenvals implies smooth MP}]
    Let $A$ be an expanding integer matrix such that for some $k\in \N$, $A^k$ is diagonalizable with integer eigenvalues. To show that $A$ has a smooth Markov partition, combine Lemmas \ref{lem:MP for diagonal} and \ref{lem:MPforsimilarmatrices} to get a linear Markov partition for $A^k$, then apply Lemma \ref{lem:MPforpowers} to get a Markov partition for $A$.
\end{proof}



\subsubsection{On the hypothesis of Theorem \ref{thm:power diag with integer eigenvals implies smooth MP}}

We point out that while Theorem \ref{thm:power diag with integer eigenvals implies smooth MP} asks that \textit{some} power of an expanding integer matrix $A$ has integer eigenvalues, the minimal such power is actually bounded by a constant depending only on the dimension. The idea of the proof is to first consider the smallest power $\nu$ for which the eigenvalues of $A^\nu$ all have integer modulus. The characteristic polynomial of $A^\nu$ then (essentially) decomposes into cyclotomic polynomials, which allow us to bound how big $n$ needs to be for $A^n$ to have integer eigenvalues.



Let us introduce some notation. Let $\mathcal{A}_d$ be the collection of $A \in M_d(\Z)$ for which there exists an $n \in \N$ such that $A^n$ has integer eigenvalues. The minimum such $n$ will be written $n_A$. Also, denote $$\mathcal{H}(d):=\{n_A: A\in \mathcal{A}_d\}$$
and $$\mathcal{L}(d):= \sup \mathcal{H}(d).$$

\begin{proposition}\label{prop:bounded_power}
    For all $d\in \N$, $\mathcal{L}(d)$ is finite.
\end{proposition}

\begin{proof}
    Let $A\in \mathcal{A}_d$ and $\lambda_j = r_je^{i\theta_j}$ with $r_j \ge 0$ denote its eigenvalues. Observe that if the characteristic polynomial $\chi_A(x)$ is a product of two polynomials $f(x), g(x) \in \Q[x]$, then $f$ and $g$ partition the roots of $\chi_A$. Hence, $n_A$ will be the least common multiple of the smallest $n_f$ and smallest $n_g$ for which $\xi^{n_f} \in \Z$ and $\omega^{n_g} \in \Z$ for all roots $\xi$ of $f$ and $\omega$ of $g$. Since $\chi_A$ consists of at most $d$ irreducible factors, it thus suffices to bound $n_A$ by a function of $d$ when $\chi_A$ is irreducible over $\Q$.
    
    Assuming $\chi_A$ is irreducible, write $k = \lambda^{n_A} \in \Z$ for some root $\lambda$ of $\chi_A$. Then, $\lambda$ is a root of $x^{n_A} - k$, so $\chi_A$ divides $x^{n_A} - k$. This means all the roots of $\chi_A$ have the same modulus, so $r_1=r_2=...=r_d=r$ for some $r\geq0$. When $r=0$ the situation is of course trivial, so we consider $r>0$. 
    
    
    Furthermore, the constant coefficient of $\chi_A$ is $\prod_{j=1}^d \lambda_j\in \Z$, so $r^d \in \N$. Thus, the smallest $\nu \in \N$ for which $r^\nu \in \N$ is at most $d$. After raising $A$ to the power $\nu$, and taking its characteristic polynomial, we obtain 
    $$
        \chi_{A^\nu}( r^\nu x) = \prod_{j=1}^d \brac{r^\nu x - r^\nu e^{i \nu\theta_j}}= r^{\nu d} \underbrace{\prod_{j=1}^d \brac{x - e^{i \nu\theta_j}}}_{\in \Q[x]}.
    $$
    Let $\varphi$ denote Euler's totient function and $\Phi_m$ the $m^{\text{th}}$ cyclotomic polynomial, so that $\deg(\Phi_m) = \varphi(m)$. By assumption, $e^{in_A\theta_j} = \pm 1$, so $\theta_j \in \pi \Q$. That is, all the roots of $\chi_{A^\nu}( r^\nu x)$ are roots of unity. Hence, $\chi_{A^\nu}( r^\nu x)/r^{\nu d}$ is a product of cyclotomic polynomials:
    $$
        \chi_{A^\nu}( r^\nu x) = r^{\nu d} \prod_{j=1}^d \brac{x - e^{i \nu\theta_j}} = r^{\nu d} \prod_{J \in \mJ} \Phi_{m_J}(x)
    $$
    for some indexing set $\set{m_J: J \in \mJ} \subset \N$. The smallest natural $n_J$ such that $\xi^{n_J} =\pm 1$ for all roots $\xi$ of $\Phi_{m_J}$ is
    $$
        n_J = \min \Set{t \in \N: \frac{2\pi t}{m_J} = 0 \mod \pi} = \frac{m_J}{\gcd(2,m_J)} \le m_J.
    $$
    Additionally, $\deg(\Phi_{m_J}) = \varphi(m_J) \le d$ since $\deg(\chi_{A^\nu}) = d$. It is well known that $\lim_{t \to \infty} \varphi(t) = \infty$, so there are only finitely many cyclotomic polynomials with degree at most $d$:
    $$
         m_J \le \max_{1 \le t \le d} \max \varphi^{-1}(t) < \infty
    $$
    Thus, we observe that
    $$
        n_A = \nu \, \lcm\set{n_J : J \in \mJ} \le d \prod_{J} n_J \le d \prod_{J} m_J \le d \brac{\max_{1 \le t \le d} \max \varphi^{-1}(t)}^d < \infty,
    $$
    as desired.
\end{proof}
After raising our matrix to the power $\nu$, the question of interest reduces to ``\textit{What possible products of cyclotomic polynomials can have total degree $d$?'}' This question is naturally related to finite order elements of the group $\operatorname{GL}_d(\Q)$. Indeed, a manual computation shows that in small dimensions 

$$
    \mathcal{H}(d)=\big\{n\in \N:\exists B\in \operatorname{GL}_d(\Q) \text{ such that } n=\inf \{k:B^k=I_d\}\big\},
$$
where $I_d$ denotes the $d\times d$ identity matrix. This latter question is well-studied, and one can use work in this direction to get a sharper bound on $\mathcal{L}(d)$, see e.g. \cite{KP, K}.

When $d=2$, the approach from the proof of Proposition \ref{prop:bounded_power} shows that $\mathcal{H}(2)= \set{1,2,3,4,6}$. One consequence of this is that in order to check if a $2\times 2$ expanding matrix $A$ has a linear Markov partition, it suffices to check if $A^{12}$ is diagonalizable with integer eigenvalues, since $12 = \lcm \set{1,2,3,4,6}$. Indeed, in dimension $d$, some power of $A$ has integer eigenvalues if and only if $A^K$ does, where $K=\lcm\mathcal{H}(d)$. This gives a conceptually simple, albeit computationally demanding, way of checking whether or not a given matrix admits a linear Markov partition. 

Showing that the values in $\set{1,2,3,4,6}$ are attainable is done by construction. Let $A\in \mathcal{A}_2$ with eigenvalues $\lambda_1,\lambda_2$ and let $\nu$ be the smallest natural such that $\abs{\lambda_j}^\nu \in \N$. Following the proof of Proposition \ref{prop:bounded_power}, we see that $\nu \in \set{1,2}$. If $\chi_{A^\nu}$ is reducible, it must be the product of two linear factors, so $A^\nu$ has integer eigenvalues. If $\chi_{A^\nu}$ is irreducible, then $A$ has two conjugate complex eigenvalues, so
$$
    \chi_A(x) = (x-re^{i\theta})(x-re^{-i\theta}) = x^2 - 2r\cos(\theta) x + r^2\in \Z[x].
$$
Hence, $2r\cos(\theta) \in \Z$. We also know that $e^{i\nu \theta}$ is a primitive $m^\text{th}$ root of unity for some $m$ such that $\deg\Phi_m = 2$, so $m \in \varphi^{-1}(2) = \set{3,4,6}$. 

If $n_A$ is, say, $4$, then $m = 4$ and $\nu = 2$. Thus, $\nu \theta = 2\pi/4 \mod 2\pi$, so $\theta \in \pi/4 + \set{0,\pi}$, meaning $\cos(\theta) = \pm\sqrt2/2$. This implies that
$$
    \chi_A(x) = x^2 \mp \sqrt 2 r x + r^2,
$$
so $r = k \sqrt 2$ for some $k \in \N$. The conclusion is that the rational cannonical form for $A$ is
$$
    \mtr{0&-2k^2\\ 1 & \pm2k}.
$$
The rational cannonical forms for $A$ such that $n_A \in \set{3,6}$ can be found similarly.

\section{Non-smoothness in dimension 2} \label{Section:NonSmooth}




In this section, we will prove Theorem \ref{thm:nonsmoothness in dimension 2} by handling each case separately. We will need the following well-known lemma from ergodic theory, which is a consequence of the Kronecker-Weyl Theorem. For a vector $\vec v\in \R^d$ and $S\subset \R$, we will write $S\vec v$ to mean the set $\{t\vec v :t\in S\}$. Denote also $\R^+:=[0,\infty)$.



\begin{lemma}\label{lem-irrationalslopemeansuniformlydense}
    Let $\vec v \in\R^2$ be a vector with irrational slope. Then for any $\varepsilon>0$, there exists an $N>0$ such that $p([0,N]\vec v )$ is $\varepsilon$-dense in $\T^2$. Moreover, for any $x_0\in\T^2$, $x_0+p([0,N]\vec v)$ is also $\varepsilon$-dense. 
\end{lemma}

\begin{proof}
    Let $\varepsilon > 0$. Since $\T^2\subset \bigcup_{x\in p(\R^+\vec{v})}B(x,\varepsilon)$, by compactness, we have finitely many $x_1,...,x_n\in p(\R^+\vec{v})$ such that $\T^2\subset \bigcup_{i=1}^nB(x_i,\varepsilon)$. But $p$ is a bijection on $\R^+\vec{v},$ so for each $1\leq i\leq n$ there exists a unique $t_i\in \R^+$ such that $x_i=p(t_i\vec{v})$. Thus, setting $N=\max_{1\leq i \leq n}\{t_i\}$,  $p([0,N]\vec{v})$ is $\varepsilon$-dense in $\T^2$. Finally, $\varepsilon$-denseness of translations is immediate.
\end{proof}

We now move to the proofs of the three cases of Theorem \ref{thm:nonsmoothness in dimension 2}. Throughout, let $A$ be a $2\times 2$ expanding integer matrix. We will be working with curve parameterizations $\gamma:[0,1]\to \R^2$ and to avoid notation that is too cumbersome, we occasionally allow ourselves write $\gamma$ to mean the set $\gamma([0,1])$. 

\subsection{Real eigenvalues}
\label{subsection: Real Eigenvals}


Suppose $\lambda_1,\lambda_2\in \R\setminus\Q$ are the eigenvalues of $A$ arranged such that $|\lambda_1|>|\lambda_2|$. Take the associated eigenvectors $\vec v_1, \vec v_2$ to be unit vectors. Recall that our objective is to show that given $\M$ a Markov partition for $f=f_A$, $\partial\M$ cannot contain a $C^1$ arc, which we will do by showing that existence of such a curve would imply that $\partial \M$ is dense, contradicting Lemma \ref{lem:boundary not dense}.

Let us suppose towards this contradiction that such a curve $C$ exists. Let $\gamma:[0,1]\to \R^2$, $x_0 \in \R^2$ so that $\gamma(0) = 0$ and $\gamma(t) + x_0$ is a parametrization of a lift of $C$ with nowhere zero velocity. Let $\proj_{\vec v_1}: \R^2 \to \R \vec v_1$ denote the projection onto $\R \vec v_1$ along $\R \vec v_2$ and define $\proj_{\vec v_2}$ analogously. Write $\gamma(t) = \gamma_1(t)\vec v_1 + \gamma_2(t) \vec v_2$ so that $\proj_{\vec v_1}(\gamma(t)) = \gamma_1(t) \vec v_1$.

Suppose first that on some subinterval of $[0,1]$ with positive length, $\gamma^\prime(t)$ is parallel to $\vec v_2$. Then $\gamma(t)$ is linear and parallel to $\vec v_2$ on this interval. Since $\vec v_2$ has irrational slope,  $$\partial\M=\bigcup_{k=0}^\infty f^k(\partial\M)\supset \bigcup_{k=0}^\infty f^k\left( \vphantom{\prod} p(x_0) +  p\gamma([0,1])\right)$$ is dense in $\T^2$, a contradiction. 

We may thus suppose otherwise, and without loss of generality adjust our interval so that $\gamma_1^\prime$ is non-vanishing. Additionally, we may traverse $\gamma$ so that $\gamma^\prime_1$, hence also $\gamma_1$, is positive. By compactness of the unit interval, there exist $m, M > 0$ such that for all $t\in [0,1]$,
$$
    m \leq \gamma_1^\prime(t), \text{ and } \abs{\gamma_2^\prime(t)} \leq M.
$$
Letting $\|\cdot\|$ denote the Euclidean norm on $\R^2$, this means
\begin{align*}
    \norm{\proj_{\vec v_1}A^{2n}\gamma(t)} & =\lambda_1^{2n} \gamma_1(t) \\
    & = \lambda_1^{2n} \int_0^{t} \gamma_1^\prime(s)ds\\
    & \geq \lambda_{1}^{2n} t m
\end{align*}
At the same time,
\begin{align*}
    \norm{\proj_{\vec v_2}A^{2n}\gamma(t)} &\le \lambda_2^{2n}\int_0^{t} \abs{\gamma_2^\prime(s)}ds\\
    &\leq \lambda_2^{2n} t M.
\end{align*}


Let now $\varepsilon > 0$. By Lemma \ref{lem-irrationalslopemeansuniformlydense}, there exists an $N > 0$ such that for all $y_0 \in \T^2$, $y_0 + p([0,N]\vec v_1)$ is $\varepsilon/2$-dense in $\T^2$. Set $R = N/m$. Since $\abs{\lambda_1} > \abs{\lambda_2}$, there exists an $n\in \N$ such that for all $0 \le r \le R$, 
\begin{align*}
    \norm{A^{2n}\gamma(r\lambda_1^{-2n}) - \proj_{\vec v_1}A^{2n}\gamma(r\lambda_1^{-2n})} & = \norm{\proj_{\vec v_2}A^{2n}\gamma(r\lambda_1^{-2n})}\\
    &\leq \brac{\frac{\lambda_2}{\lambda_1}}^{2n} rM \\
    &\leq \brac{\frac{\lambda_2}{\lambda_1}}^{2n} RM \\
    &< \varepsilon/2
\end{align*}
Additionally,
$$
    \norm{\proj_{\vec v_1}A^{2n}\gamma(R \lambda_1^{-2n})} \geq Rm = N \quad \text{and} \quad \proj_{\vec v_1}A^{2n}\gamma(0)=0,
$$
so for all $0 \leq \ell \leq N$, there exists an $r \in [0,R]$ such that $\proj_{\vec v_1}A^{2n}\gamma(r\lambda_1^{-2n}) = \ell \vec v_1$ by the Intermediate Value Theorem. In other words, $[0,N]\vec v_1$ is contained in an $\varepsilon/2$-neighbourhood of $A^{2n} \gamma$. Projecting down to $\T^2$, we see that
$$
    \partial \mM \supset f^{2n}C = pA^{2n}(x_0+\gamma([0,1]))
$$
is $\varepsilon$-dense in $\T^2$ by the $\varepsilon/2$-density of $p(A^{2n}x_0) + p([0,N]\vec v_1)$. Our choice of $\varepsilon$ was arbitrary, so $\partial \mM$ is dense in $\T^2$: we have reached the desired contradiction.


\begin{figure}
    \centering
    \begin{overpic}[width=0.4\textwidth]{620_gamma_iterates_shortv4.png}
        \put(75,225){{\Large $\vec v_2$}}
        \put(910,260){{\Large $\vec v_1$}}
        \put(250,250){{\Large $\gamma$}}
        \put(320,380){{\Large $A\gamma$}}
        \put(510,470){{\Large $A^2\gamma$}}
        \put(700,465){{\Large $A^3\gamma$}}
    \end{overpic}
    \begin{overpic}[width=0.4\textwidth]{620_gamma_iterates_proj_map.png}
        \put(195,315){{\Large $A\gamma(t)$}}
        \put(-120,10){{\large $\proj_{\vec v_2} A\gamma(t)$}}
        \put(650,100){{\large $\proj_{\vec v_1} A\gamma(t)$}}
    \end{overpic}
    \caption{ Iterates of $\gamma$ under a matrix $A$ with real, irrational eigenvalues of different modulus and the projection maps onto the respective eigenlines.}
    \label{fig:irrational_eigenvalues}
\end{figure}

\subsection{Complex eigenvalues}

Suppose $A$ has complex eigenvalues $\lambda e^{i\phi},\lambda e^{-i\phi}$ such that $\phi\notin \pi\Q$ and $\lambda > 1$. It is well known that up to a linear change of coordinates, $A$ acts on the plane by scaling by $\lambda$ and rotating by $\phi$. Let $T \in GL_2(\R)$ be such that $A = T\lambda R_{\phi} T^{-1}$ where $R_\phi$ is a rotation matrix.

Let $\M$ be a Markov partition for $f_A$. Following Definition \ref{def:smoothness MP}, we suppose towards a contradiction that a continuous injection $\tilde\gamma:[0,1]\to \T^2$ exists with $\tilde\gamma'(t)\neq 0$ for some $t$. Without loss of generality, we can parametrize this curve as follows. Let $x_0 \in \R^2$, $\gamma:[0,1] \to \R^2$ continuous and injective such that $\gamma(0) = 0$, $\gamma^\prime(0) \neq 0$, and $T(x_0 + \gamma(t))$ parametrizes a lift of $\tilde\gamma$.

Let $\vec{v}$ be an arbitrary vector with irrational slope. Similarly to Section \ref{subsection: Real Eigenvals}, it suffices by Lemma \ref{lem-irrationalslopemeansuniformlydense} to show that for all $N> 0$ and $\varepsilon > 0$, there exists an $n$ such that $[0,N] \vec v$ is contained in an $\varepsilon$-neighbourhood of
$$
    A^n T\gamma([0,1]) = T \lambda^n R^n_\phi\gamma([0,1]).
$$
This would imply that $\partial\M$ is dense in $\T^2$, contradicting Lemma \ref{lem:boundary not dense}. However, $T$ is a homeomorphism and $\lambda^n R^n_\phi\gamma([0,1])$ is compact. So if $\mathcal{U}$ is an $\varepsilon$-neighbourhood of $T\lambda^n R^n_\phi\gamma([0,1])$, then it is not hard to show that there exists an $\varepsilon^\prime$ and an $\varepsilon^\prime$-neighbourhood $V$ of $\lambda^n R^n_\phi\gamma([0,1])$ satisfying $T(V) \subset \mathcal{U}$. That is to say, it suffices to show that $[0,N] T^{-1} \vec v$ is contained in an $\varepsilon$-neighbourhood of $\lambda^n R^n_\phi\gamma([0,1])$.

Fix $\varepsilon, N > 0$. Let $\vec w = T^{-1} \vec v$, rescaled to have unit length, and for non-zero $\vec a, \vec b \in \R^2$, let $\theta(\vec a, \vec b) \in (-\pi,\pi]$ denote the signed angle between $\vec a$ and $\vec b$. Since $\phi \notin \pi\Q$, we can pick a sequence $(n_k)_k$ such that $$ \lim_{k\to\infty}\theta( R^{n_k}_\phi\gamma^\prime(0), \vec w) = 0.$$ As $R_\phi$ is conformal, $$\theta(R_\phi^{n_k} \gamma(t), R_\phi^{n_k} \gamma^\prime(0)) = \theta(\gamma(t), \gamma^\prime(0))$$ for all $k$, and by the definition of $\gamma^\prime(0)$, $$ \lim_{t\to 0}\theta(\gamma(t), \gamma^\prime(0)) = 0.$$ 

Let $\delta<\pi/2$ small enough so that $\tan\delta<\varepsilon/N$. Then, there exist $\eta,K > 0$ such that whenever $t\in (0,\eta]$ and $k > K$, the following holds:
\begin{align*}
    \abs{\theta(R_\phi^{n_k} \gamma(t), \vec w)} & \le \abs{\theta(R_\phi^{n_k} \gamma(t), R_\phi^{n_k}\gamma^\prime(0))} + \abs{ \theta(R_\phi^{n_k}\gamma^\prime(0),  \vec w)}\\
    & = \abs{\theta( \gamma(t), \gamma^\prime(0))} + \abs{\theta(R_\phi^{n_k}\gamma^\prime(0),  \vec w)}\\
    & < \delta.
\end{align*}
In particular, our choice of $\delta$ means we have $\vec w \cdot R_\phi^{n_k} \gamma(t) > 0$. Fix $k > K$ so that
$$
    \vec w \cdot \lambda^{n_k} R_\phi^{n_k} \gamma(\eta) > N.
$$

Let $p_{\vec{w}}:\vec {a} \mapsto (\vec {a} \cdot \vec {w}) \vec {w}$ denote the orthogonal projection onto $\R \vec w$. As in Section \ref{subsection: Real Eigenvals}, the Intermediate Value Theorem now tells us that for every $0 \le \ell \le N$ there exists a $t\in [0,\eta]$ such that $$p_{\vec{w}}\lambda^{n_k}R_{\phi}^{n_k}\gamma(t)=\ell\vec{w}.$$ Moreover, for $t\neq 0$,
\begin{align*}
    \norm{\lambda^{n_k} R_\phi^{n_k} \gamma(t) -p_{\vec w} \lambda^{n_k}  R_\phi^{n_k} \gamma(t)} & =  \norm{p_{\vec w}\lambda^{n_k} R_\phi^{n_k} \gamma(t)} \abs{\tan\theta(R_\phi^{n_k} \gamma(t), \vec w)}\\
    & < N \tan \delta \\
    & < \varepsilon
\end{align*}
and the same inequality is trivially satisfied when $t=0$. 

We conclude that $[0,N]\vec w$ is contained in an $\varepsilon$-neighbourhood of $\lambda^{n_k}R_\phi^{n_k}\gamma([0,\eta]) \subset \lambda^{n_k}R_\phi^{n_k}\gamma([0,1])$, as desired.


\begin{figure}
    \centering
    \includegraphics[width=0.5\linewidth]{620_gamma_rotation_iterates.png}
    \caption{Iterates of $\gamma$ under a matrix $A$ with complex eigenvalues whose arguments do not lie in $\pi\Q $. $L=\R^+\vec{w}$ and $n_k$ are chosen so that $A^{n_k} \gamma^\prime(0)$ approaches the direction of $L$ counterclockwise.}
    \label{fig:complex_eigenvalues}
\end{figure}

\subsection{Non-diagonalizable matrices}
If $A$ is not diagonalizable, then $A$ is similar over $\R$ to some Jordan block $J$. It is straightforward to show that the eigenvalue $k$ of $A$ is an integer, so $A$ is similar over $\Q$ to $J \in M_2(\Z)$. Note that by Proposition \ref{prop:MPs give Markov tilings} and Lemma \ref{lem:MPforsimilarmatrices}, it is sufficient to show that $J$ does not admit a smooth Markov partition. It is also straightforward to show that the set
$$
    D = \Set{\alpha \in \R : \Set{k^{2n} \alpha \ \text{ mod } 1}_{n\in\N} \text{ is dense in }[0,1]}
$$
is dense in $\R$. 

Suppose for the sake of contradiction that $J$ has a smooth Markov partition $\M$. We will once again show that smoothness of $\partial\M$ contradicts Lemma \ref{lem:boundary not dense}. However, unlike the previous cases, we will do so by finding a particular curve contained in $\partial \M$ with some additional properties.

Since $J$ is expanding, $\abs{k} \ge 2$. Hence, $\partial \M$ cannot purely consist of a union of lines parallel to the $x$-axis. Such a partition consists of horizontal strips, which would get wrapped $k$-fold around the torus horizontally after applying $J$. Thus, one would be able to find two points in the interior of some rectangle that have the same image in the interior of another rectangle, contradicting the Markov property (MP).


Therefore, there exists an arc $C \subset \partial\M$ with a lift parametrized by a $C^1$ map $\gamma:[0,1] \to \R^2$ such that the vertical component of $\gamma^\prime$ is non-zero. For a vector $\vec{v}\in \R^2$ written in the standard basis as $\vec{v}=x\vec e_1+y\vec e_2 $, let $$p_1(\vec{v})=x,\quad p_2(\vec{v})=y.$$ Since the vertical component of $\gamma^\prime$ is non-vanishing, $p_2 \gamma$ contains an open interval, so we may assume $p_2\gamma(0) \in D$ by the density of $D$. We also traverse $\gamma$ so that $p_2\gamma'$ has the same sign as $k$. 

Observe that

$$
    J^n = \brac{\mtr{k & 0 \\ 0 & k} + \mtr{0 & 1 \\ 0 & 0}}^n = \sum_{l=0}^n {n \choose l} \mtr{k & 0 \\ 0 & k}^{n-l} \mtr{0 & 1 \\ 0 & 0}^l = \mtr{k^n & nk^{n-1} \\ 0 & k^n},
$$
so for any $a,b\in \R$, $$J^{2n}\begin{pmatrix} a \\ b \end{pmatrix} = \begin{pmatrix}
    k^{2n}a+2nk^{2n-1}b \\
    k^{2n}b
\end{pmatrix}.$$ 

We use a similar convergence argument as in Section \ref{subsection: Real Eigenvals}, but here the $x$-axis plays the role of the dominant eigendirection. By continuity and compactness of $[0,1]$ there exist constants $M_1,M_2>0$, and $m_2\neq 0$ with the same sign as $k$ and $p_2\gamma'$ such that $$|p_1\gamma'(t)|\leq M_1, \quad 0< |m_2|\leq |p_2\gamma'(t)|\leq M_2.$$ Let $\varepsilon>0$ and $x\in[0,1)^2$ be arbitrary. Since $p_2J^{2n}\gamma(0)=k^{2n}p_2\gamma(0)$ and $p_2\gamma(0)\in D$, there exists an $$n>\frac{1+\varepsilon\frac{M_1}{M_2}}{2\varepsilon\frac{m_2}{kM_2}}$$ for which $|p_2(x)-(p_2J^{2n}\gamma(0)\ \text{ mod } 1)|<\varepsilon.$

Let $$z=p_1(x)-p_1J^{2n}\gamma(0)\ \text{ mod } 1$$
and denote $$\eta_n:=\frac{\varepsilon}{k^{2n}M_2}.$$

Then by our choice of $n$, \begin{align*}
    p_1J^{2n}(\gamma(\eta_n)-\gamma(0))&=\int_0^{\eta_n} p_1J^{2n}\gamma'(s)ds \\
    &=\int_0^{\eta_n} k^{2n}p_1\gamma'(s)+2nk^{2n-1}p_2\gamma'(s)ds \\
    &\geq -\eta_nk^{2n}M_1 + \eta_n2nk^{2n-1}m_2 \\
    &= -\varepsilon\frac{M_1}{M_2}+2\varepsilon n\frac{m_2}{kM_2} > 1
\end{align*}
Thus, by continuity and the Intermediate Value Theorem, there exists a $t\in[0,\eta_n]$ such that $$p_1J^{2n}(\gamma(t)-\gamma(0))=z.$$ Moreover, \begin{align*}
    \big|p_2J^{2n}(\gamma(t)-\gamma(0))\big|&\leq \int_0^t\big|p_2J^{2n}\gamma'(s)\big|ds \\
    &= \int_0^t \big|k^{2n}p_2\gamma'(s)\big|ds \\
    &\leq tk^{2n}M_2 \\
    &\leq \eta_nk^{2n}M_2 \\
    &= \varepsilon
\end{align*}

By construction, there is some $\Delta\in \Z^2$ for which \begin{align*}
    \big\|x-J^{2n}\gamma(t)-\Delta\big\| &\leq \big|p_1(x-J^{2n}\gamma(t)-\Delta)\big| + \big|p_2(x-J^{2n}\gamma(t)-\Delta)\big| \\
    &\leq \big|p_1(x)-z-p_1J^{2n}\gamma(0)-p_1(\Delta)\big| \\
    &\quad\;\; + \big|p_2(x)-p_2J^{2n}\gamma(0)-p_2(\Delta)\big| + \big|p_2J^{2n}(\gamma(t)-\gamma(0))\big| \\
    &<0+\varepsilon+\varepsilon=2\varepsilon
\end{align*} 

Projecting down to $\T^2$, this means $d\big(p(x),f^{2n}p\gamma(t)\big)<2\varepsilon$. Our choices for $x$ and $\varepsilon$ were arbitrary, so we conclude that 
$$
    \partial\M=\bigcup_{k=0}^\infty f^k(\partial \M) \supset \bigcup_{k=0}^\infty f^{2k} C 
$$
is dense in $\T^2$. This contradicts Lemma \ref{lem:boundary not dense}, and concludes the proof of Theorem \ref{thm:nonsmoothness in dimension 2}.\qed 
\\

\begin{figure}
    \centering
    \includegraphics[width=0.5\linewidth]{620_gamma_jordan_iterates.png}
    \caption{Iterates of $\gamma$ under a Jordan matrix $A$. The $x$-axis is scaled down and $n_k$ are some large integers.}
    \label{fig:jordan_matrix_curves}
\end{figure}


The contradiction in the proof above only arises from the existence of a $C^1$ curve with vertical change contained in $\partial \M$, which does not rule out linear behaviour in the horizontal direction. Indeed, in the case that $$A=\begin{pmatrix}
        2 & 1 \\ 0 & 2
\end{pmatrix}$$ a construction due to Bedford \cite{Be} yields the Markov partition in Figure \ref{fig:hybrid MP}, for which the boundary along the $x$-axis is linear, but the vertical\footnote{For lack of better terminology.} component is fractal. In fact, there are straight line segments at every dyadic rational.

Figure \ref{fig:hybrid MP} is generated by computer code executing the first few steps of the iterative procedure in \cite{Be}, using anticlockwise perturbations of integer-valued vectors. The dashed green line indicates a fundamental domain for $\T^2$, so the displayed partition consists of four rectangles. 

For general Jordan matrices $$\begin{pmatrix}
    k & 1 \\ 0 & k
\end{pmatrix} \;\text{ with }\; k\in \Z\setminus \{-1,0,1\},$$ the procedure yields a similar partition with $k^2$ rectangles. In other words, Jordan matrices allow for Markov partitions with hybrid behaviour. Using this construction, we can build a hybrid Markov partition for any non-diagonalizable $2\times 2$ matrix $A$.

\begin{figure}
    
    \centering

    \begin{circuitikz}
        \def\Aa(#1,#2){(0.03125*#1-0.078125*#2,0.03125*#2)}
        

        
        \draw[short] \Aa(80,32) -- \Aa(81,32);
        

        \def\Aa(#1,#2){(0.03125*#1-0.078125*#2+1,0.03125*#2)}
        
        \draw[short] \Aa(0,0) -- \Aa(0,1);
        \draw[short] \Aa(0,1) -- \Aa(0,2);
        \draw[short] \Aa(0,2) -- \Aa(1,2);
        \draw[short] \Aa(1,2) -- \Aa(1,3);
        \draw[short] \Aa(1,3) -- \Aa(1,4);
        \draw[short] \Aa(1,4) -- \Aa(2,4);
        \draw[short] \Aa(2,4) -- \Aa(3,4);
        \draw[short] \Aa(3,4) -- \Aa(4,4);
        \draw[short] \Aa(4,4) -- \Aa(4,5);
        \draw[short] \Aa(4,5) -- \Aa(4,6);
        \draw[short] \Aa(4,6) -- \Aa(5,6);
        \draw[short] \Aa(5,6) -- \Aa(5,7);
        \draw[short] \Aa(5,7) -- \Aa(5,8);
        \draw[short] \Aa(5,8) -- \Aa(6,8);
        \draw[short] \Aa(6,8) -- \Aa(7,8);
        \draw[short] \Aa(7,8) -- \Aa(8,8);
        \draw[short] \Aa(8,8) -- \Aa(9,8);
        \draw[short] \Aa(9,8) -- \Aa(10,8);
        \draw[short] \Aa(10,8) -- \Aa(11,8);
        \draw[short] \Aa(11,8) -- \Aa(12,8);
        \draw[short] \Aa(12,8) -- \Aa(12,9);
        \draw[short] \Aa(12,9) -- \Aa(12,10);
        \draw[short] \Aa(12,10) -- \Aa(13,10);
        \draw[short] \Aa(13,10) -- \Aa(13,11);
        \draw[short] \Aa(13,11) -- \Aa(13,12);
        \draw[short] \Aa(13,12) -- \Aa(14,12);
        \draw[short] \Aa(14,12) -- \Aa(15,12);
        \draw[short] \Aa(15,12) -- \Aa(16,12);
        \draw[short] \Aa(16,12) -- \Aa(16,13);
        \draw[short] \Aa(16,13) -- \Aa(16,14);
        \draw[short] \Aa(16,14) -- \Aa(17,14);
        \draw[short] \Aa(17,14) -- \Aa(17,15);
        \draw[short] \Aa(17,15) -- \Aa(17,16);
        \draw[short] \Aa(17,16) -- \Aa(18,16);
        \draw[short] \Aa(18,16) -- \Aa(19,16);
        \draw[short] \Aa(19,16) -- \Aa(20,16);
        \draw[short] \Aa(20,16) -- \Aa(21,16);
        \draw[short] \Aa(21,16) -- \Aa(22,16);
        \draw[short] \Aa(22,16) -- \Aa(23,16);
        \draw[short] \Aa(23,16) -- \Aa(24,16);
        \draw[short] \Aa(24,16) -- \Aa(25,16);
        \draw[short] \Aa(25,16) -- \Aa(26,16);
        \draw[short] \Aa(26,16) -- \Aa(27,16);
        \draw[short] \Aa(27,16) -- \Aa(28,16);
        \draw[short] \Aa(28,16) -- \Aa(29,16);
        \draw[short] \Aa(29,16) -- \Aa(30,16);
        \draw[short] \Aa(30,16) -- \Aa(31,16);
        \draw[short] \Aa(31,16) -- \Aa(32,16);
        \draw[short] \Aa(32,16) -- \Aa(32,17);
        \draw[short] \Aa(32,17) -- \Aa(32,18);
        \draw[short] \Aa(32,18) -- \Aa(33,18);
        \draw[short] \Aa(33,18) -- \Aa(33,19);
        \draw[short] \Aa(33,19) -- \Aa(33,20);
        \draw[short] \Aa(33,20) -- \Aa(34,20);
        \draw[short] \Aa(34,20) -- \Aa(35,20);
        \draw[short] \Aa(35,20) -- \Aa(36,20);
        \draw[short] \Aa(36,20) -- \Aa(36,21);
        \draw[short] \Aa(36,21) -- \Aa(36,22);
        \draw[short] \Aa(36,22) -- \Aa(37,22);
        \draw[short] \Aa(37,22) -- \Aa(37,23);
        \draw[short] \Aa(37,23) -- \Aa(37,24);
        \draw[short] \Aa(37,24) -- \Aa(38,24);
        \draw[short] \Aa(38,24) -- \Aa(39,24);
        \draw[short] \Aa(39,24) -- \Aa(40,24);
        \draw[short] \Aa(40,24) -- \Aa(41,24);
        \draw[short] \Aa(41,24) -- \Aa(42,24);
        \draw[short] \Aa(42,24) -- \Aa(43,24);
        \draw[short] \Aa(43,24) -- \Aa(44,24);
        \draw[short] \Aa(44,24) -- \Aa(44,25);
        \draw[short] \Aa(44,25) -- \Aa(44,26);
        \draw[short] \Aa(44,26) -- \Aa(45,26);
        \draw[short] \Aa(45,26) -- \Aa(45,27);
        \draw[short] \Aa(45,27) -- \Aa(45,28);
        \draw[short] \Aa(45,28) -- \Aa(46,28);
        \draw[short] \Aa(46,28) -- \Aa(47,28);
        \draw[short] \Aa(47,28) -- \Aa(48,28);
        \draw[short] \Aa(48,28) -- \Aa(48,29);
        \draw[short] \Aa(48,29) -- \Aa(48,30);
        \draw[short] \Aa(48,30) -- \Aa(49,30);
        \draw[short] \Aa(49,30) -- \Aa(49,31);
        \draw[short] \Aa(49,31) -- \Aa(49,32);
        \draw[short] \Aa(49,32) -- \Aa(50,32);
        \draw[short] \Aa(50,32) -- \Aa(51,32);
        \draw[short] \Aa(51,32) -- \Aa(52,32);
        \draw[short] \Aa(52,32) -- \Aa(53,32);
        \draw[short] \Aa(53,32) -- \Aa(54,32);
        \draw[short] \Aa(54,32) -- \Aa(55,32);
        \draw[short] \Aa(55,32) -- \Aa(56,32);
        \draw[short] \Aa(56,32) -- \Aa(57,32);
        \draw[short] \Aa(57,32) -- \Aa(58,32);
        \draw[short] \Aa(58,32) -- \Aa(59,32);
        \draw[short] \Aa(59,32) -- \Aa(60,32);
        \draw[short] \Aa(60,32) -- \Aa(61,32);
        \draw[short] \Aa(61,32) -- \Aa(62,32);
        \draw[short] \Aa(62,32) -- \Aa(63,32);
        \draw[short] \Aa(63,32) -- \Aa(64,32);
        \draw[short] \Aa(64,32) -- \Aa(65,32);
        \draw[short] \Aa(65,32) -- \Aa(66,32);
        \draw[short] \Aa(66,32) -- \Aa(67,32);
        \draw[short] \Aa(67,32) -- \Aa(68,32);
        \draw[short] \Aa(68,32) -- \Aa(69,32);
        \draw[short] \Aa(69,32) -- \Aa(70,32);
        \draw[short] \Aa(70,32) -- \Aa(71,32);
        \draw[short] \Aa(71,32) -- \Aa(72,32);
        \draw[short] \Aa(72,32) -- \Aa(73,32);
        \draw[short] \Aa(73,32) -- \Aa(74,32);
        \draw[short] \Aa(74,32) -- \Aa(75,32);
        \draw[short] \Aa(75,32) -- \Aa(76,32);
        \draw[short] \Aa(76,32) -- \Aa(77,32);
        \draw[short] \Aa(77,32) -- \Aa(78,32);
        \draw[short] \Aa(78,32) -- \Aa(79,32);
        \draw[short] \Aa(79,32) -- \Aa(80,32);
        \draw[short] \Aa(80,32) -- \Aa(81,32);


        \def\Aa(#1,#2){(0.03125*#1-0.078125*#2+2,0.03125*#2)}
        
        \draw[short] \Aa(0,0) -- \Aa(0,1);
        \draw[short] \Aa(0,1) -- \Aa(0,2);
        \draw[short] \Aa(0,2) -- \Aa(1,2);
        \draw[short] \Aa(1,2) -- \Aa(1,3);
        \draw[short] \Aa(1,3) -- \Aa(1,4);
        \draw[short] \Aa(1,4) -- \Aa(2,4);
        \draw[short] \Aa(2,4) -- \Aa(3,4);
        \draw[short] \Aa(3,4) -- \Aa(4,4);
        \draw[short] \Aa(4,4) -- \Aa(4,5);
        \draw[short] \Aa(4,5) -- \Aa(4,6);
        \draw[short] \Aa(4,6) -- \Aa(5,6);
        \draw[short] \Aa(5,6) -- \Aa(5,7);
        \draw[short] \Aa(5,7) -- \Aa(5,8);
        \draw[short] \Aa(5,8) -- \Aa(6,8);
        \draw[short] \Aa(6,8) -- \Aa(7,8);
        \draw[short] \Aa(7,8) -- \Aa(8,8);
        \draw[short] \Aa(8,8) -- \Aa(9,8);
        \draw[short] \Aa(9,8) -- \Aa(10,8);
        \draw[short] \Aa(10,8) -- \Aa(11,8);
        \draw[short] \Aa(11,8) -- \Aa(12,8);
        \draw[short] \Aa(12,8) -- \Aa(12,9);
        \draw[short] \Aa(12,9) -- \Aa(12,10);
        \draw[short] \Aa(12,10) -- \Aa(13,10);
        \draw[short] \Aa(13,10) -- \Aa(13,11);
        \draw[short] \Aa(13,11) -- \Aa(13,12);
        \draw[short] \Aa(13,12) -- \Aa(14,12);
        \draw[short] \Aa(14,12) -- \Aa(15,12);
        \draw[short] \Aa(15,12) -- \Aa(16,12);
        \draw[short] \Aa(16,12) -- \Aa(16,13);
        \draw[short] \Aa(16,13) -- \Aa(16,14);
        \draw[short] \Aa(16,14) -- \Aa(17,14);
        \draw[short] \Aa(17,14) -- \Aa(17,15);
        \draw[short] \Aa(17,15) -- \Aa(17,16);
        \draw[short] \Aa(17,16) -- \Aa(18,16);
        \draw[short] \Aa(18,16) -- \Aa(19,16);
        \draw[short] \Aa(19,16) -- \Aa(20,16);
        \draw[short] \Aa(20,16) -- \Aa(21,16);
        \draw[short] \Aa(21,16) -- \Aa(22,16);
        \draw[short] \Aa(22,16) -- \Aa(23,16);
        \draw[short] \Aa(23,16) -- \Aa(24,16);
        \draw[short] \Aa(24,16) -- \Aa(25,16);
        \draw[short] \Aa(25,16) -- \Aa(26,16);
        \draw[short] \Aa(26,16) -- \Aa(27,16);
        \draw[short] \Aa(27,16) -- \Aa(28,16);
        \draw[short] \Aa(28,16) -- \Aa(29,16);
        \draw[short] \Aa(29,16) -- \Aa(30,16);
        \draw[short] \Aa(30,16) -- \Aa(31,16);
        \draw[short] \Aa(31,16) -- \Aa(32,16);
        \draw[short] \Aa(32,16) -- \Aa(32,17);
        \draw[short] \Aa(32,17) -- \Aa(32,18);
        \draw[short] \Aa(32,18) -- \Aa(33,18);
        \draw[short] \Aa(33,18) -- \Aa(33,19);
        \draw[short] \Aa(33,19) -- \Aa(33,20);
        \draw[short] \Aa(33,20) -- \Aa(34,20);
        \draw[short] \Aa(34,20) -- \Aa(35,20);
        \draw[short] \Aa(35,20) -- \Aa(36,20);
        \draw[short] \Aa(36,20) -- \Aa(36,21);
        \draw[short] \Aa(36,21) -- \Aa(36,22);
        \draw[short] \Aa(36,22) -- \Aa(37,22);
        \draw[short] \Aa(37,22) -- \Aa(37,23);
        \draw[short] \Aa(37,23) -- \Aa(37,24);
        \draw[short] \Aa(37,24) -- \Aa(38,24);
        \draw[short] \Aa(38,24) -- \Aa(39,24);
        \draw[short] \Aa(39,24) -- \Aa(40,24);
        \draw[short] \Aa(40,24) -- \Aa(41,24);
        \draw[short] \Aa(41,24) -- \Aa(42,24);
        \draw[short] \Aa(42,24) -- \Aa(43,24);
        \draw[short] \Aa(43,24) -- \Aa(44,24);
        \draw[short] \Aa(44,24) -- \Aa(44,25);
        \draw[short] \Aa(44,25) -- \Aa(44,26);
        \draw[short] \Aa(44,26) -- \Aa(45,26);
        \draw[short] \Aa(45,26) -- \Aa(45,27);
        \draw[short] \Aa(45,27) -- \Aa(45,28);
        \draw[short] \Aa(45,28) -- \Aa(46,28);
        \draw[short] \Aa(46,28) -- \Aa(47,28);
        \draw[short] \Aa(47,28) -- \Aa(48,28);
        \draw[short] \Aa(48,28) -- \Aa(48,29);
        \draw[short] \Aa(48,29) -- \Aa(48,30);
        \draw[short] \Aa(48,30) -- \Aa(49,30);
        \draw[short] \Aa(49,30) -- \Aa(49,31);
        \draw[short] \Aa(49,31) -- \Aa(49,32);
        \draw[short] \Aa(49,32) -- \Aa(50,32);
        \draw[short] \Aa(50,32) -- \Aa(51,32);
        \draw[short] \Aa(51,32) -- \Aa(52,32);
        \draw[short] \Aa(52,32) -- \Aa(53,32);
        \draw[short] \Aa(53,32) -- \Aa(54,32);
        \draw[short] \Aa(54,32) -- \Aa(55,32);
        \draw[short] \Aa(55,32) -- \Aa(56,32);
        \draw[short] \Aa(56,32) -- \Aa(57,32);
        \draw[short] \Aa(57,32) -- \Aa(58,32);
        \draw[short] \Aa(58,32) -- \Aa(59,32);
        \draw[short] \Aa(59,32) -- \Aa(60,32);
        \draw[short] \Aa(60,32) -- \Aa(61,32);
        \draw[short] \Aa(61,32) -- \Aa(62,32);
        \draw[short] \Aa(62,32) -- \Aa(63,32);
        \draw[short] \Aa(63,32) -- \Aa(64,32);
        \draw[short] \Aa(64,32) -- \Aa(65,32);
        \draw[short] \Aa(65,32) -- \Aa(66,32);
        \draw[short] \Aa(66,32) -- \Aa(67,32);
        \draw[short] \Aa(67,32) -- \Aa(68,32);
        \draw[short] \Aa(68,32) -- \Aa(69,32);
        \draw[short] \Aa(69,32) -- \Aa(70,32);
        \draw[short] \Aa(70,32) -- \Aa(71,32);
        \draw[short] \Aa(71,32) -- \Aa(72,32);
        \draw[short] \Aa(72,32) -- \Aa(73,32);
        \draw[short] \Aa(73,32) -- \Aa(74,32);
        \draw[short] \Aa(74,32) -- \Aa(75,32);
        \draw[short] \Aa(75,32) -- \Aa(76,32);
        \draw[short] \Aa(76,32) -- \Aa(77,32);
        \draw[short] \Aa(77,32) -- \Aa(78,32);
        \draw[short] \Aa(78,32) -- \Aa(79,32);
        \draw[short] \Aa(79,32) -- \Aa(80,32);
        

        \def\Aa(#1,#2){(0.03125*#1-0.078125*#2,0.03125*#2+1)}
        
        
        \draw[short] \Aa(80,32) -- \Aa(81,32);

        \def\Aa(#1,#2){(0.03125*#1-0.078125*#2+1,0.03125*#2+1)}
        
        \draw[short] \Aa(0,0) -- \Aa(0,1);
        \draw[short] \Aa(0,1) -- \Aa(0,2);
        \draw[short] \Aa(0,2) -- \Aa(1,2);
        \draw[short] \Aa(1,2) -- \Aa(1,3);
        \draw[short] \Aa(1,3) -- \Aa(1,4);
        \draw[short] \Aa(1,4) -- \Aa(2,4);
        \draw[short] \Aa(2,4) -- \Aa(3,4);
        \draw[short] \Aa(3,4) -- \Aa(4,4);
        \draw[short] \Aa(4,4) -- \Aa(4,5);
        \draw[short] \Aa(4,5) -- \Aa(4,6);
        \draw[short] \Aa(4,6) -- \Aa(5,6);
        \draw[short] \Aa(5,6) -- \Aa(5,7);
        \draw[short] \Aa(5,7) -- \Aa(5,8);
        \draw[short] \Aa(5,8) -- \Aa(6,8);
        \draw[short] \Aa(6,8) -- \Aa(7,8);
        \draw[short] \Aa(7,8) -- \Aa(8,8);
        \draw[short] \Aa(8,8) -- \Aa(9,8);
        \draw[short] \Aa(9,8) -- \Aa(10,8);
        \draw[short] \Aa(10,8) -- \Aa(11,8);
        \draw[short] \Aa(11,8) -- \Aa(12,8);
        \draw[short] \Aa(12,8) -- \Aa(12,9);
        \draw[short] \Aa(12,9) -- \Aa(12,10);
        \draw[short] \Aa(12,10) -- \Aa(13,10);
        \draw[short] \Aa(13,10) -- \Aa(13,11);
        \draw[short] \Aa(13,11) -- \Aa(13,12);
        \draw[short] \Aa(13,12) -- \Aa(14,12);
        \draw[short] \Aa(14,12) -- \Aa(15,12);
        \draw[short] \Aa(15,12) -- \Aa(16,12);
        \draw[short] \Aa(16,12) -- \Aa(16,13);
        \draw[short] \Aa(16,13) -- \Aa(16,14);
        \draw[short] \Aa(16,14) -- \Aa(17,14);
        \draw[short] \Aa(17,14) -- \Aa(17,15);
        \draw[short] \Aa(17,15) -- \Aa(17,16);
        \draw[short] \Aa(17,16) -- \Aa(18,16);
        \draw[short] \Aa(18,16) -- \Aa(19,16);
        \draw[short] \Aa(19,16) -- \Aa(20,16);
        \draw[short] \Aa(20,16) -- \Aa(21,16);
        \draw[short] \Aa(21,16) -- \Aa(22,16);
        \draw[short] \Aa(22,16) -- \Aa(23,16);
        \draw[short] \Aa(23,16) -- \Aa(24,16);
        \draw[short] \Aa(24,16) -- \Aa(25,16);
        \draw[short] \Aa(25,16) -- \Aa(26,16);
        \draw[short] \Aa(26,16) -- \Aa(27,16);
        \draw[short] \Aa(27,16) -- \Aa(28,16);
        \draw[short] \Aa(28,16) -- \Aa(29,16);
        \draw[short] \Aa(29,16) -- \Aa(30,16);
        \draw[short] \Aa(30,16) -- \Aa(31,16);
        \draw[short] \Aa(31,16) -- \Aa(32,16);
        \draw[short] \Aa(32,16) -- \Aa(32,17);
        \draw[short] \Aa(32,17) -- \Aa(32,18);
        \draw[short] \Aa(32,18) -- \Aa(33,18);
        \draw[short] \Aa(33,18) -- \Aa(33,19);
        \draw[short] \Aa(33,19) -- \Aa(33,20);
        \draw[short] \Aa(33,20) -- \Aa(34,20);
        \draw[short] \Aa(34,20) -- \Aa(35,20);
        \draw[short] \Aa(35,20) -- \Aa(36,20);
        \draw[short] \Aa(36,20) -- \Aa(36,21);
        \draw[short] \Aa(36,21) -- \Aa(36,22);
        \draw[short] \Aa(36,22) -- \Aa(37,22);
        \draw[short] \Aa(37,22) -- \Aa(37,23);
        \draw[short] \Aa(37,23) -- \Aa(37,24);
        \draw[short] \Aa(37,24) -- \Aa(38,24);
        \draw[short] \Aa(38,24) -- \Aa(39,24);
        \draw[short] \Aa(39,24) -- \Aa(40,24);
        \draw[short] \Aa(40,24) -- \Aa(41,24);
        \draw[short] \Aa(41,24) -- \Aa(42,24);
        \draw[short] \Aa(42,24) -- \Aa(43,24);
        \draw[short] \Aa(43,24) -- \Aa(44,24);
        \draw[short] \Aa(44,24) -- \Aa(44,25);
        \draw[short] \Aa(44,25) -- \Aa(44,26);
        \draw[short] \Aa(44,26) -- \Aa(45,26);
        \draw[short] \Aa(45,26) -- \Aa(45,27);
        \draw[short] \Aa(45,27) -- \Aa(45,28);
        \draw[short] \Aa(45,28) -- \Aa(46,28);
        \draw[short] \Aa(46,28) -- \Aa(47,28);
        \draw[short] \Aa(47,28) -- \Aa(48,28);
        \draw[short] \Aa(48,28) -- \Aa(48,29);
        \draw[short] \Aa(48,29) -- \Aa(48,30);
        \draw[short] \Aa(48,30) -- \Aa(49,30);
        \draw[short] \Aa(49,30) -- \Aa(49,31);
        \draw[short] \Aa(49,31) -- \Aa(49,32);
        \draw[short] \Aa(49,32) -- \Aa(50,32);
        \draw[short] \Aa(50,32) -- \Aa(51,32);
        \draw[short] \Aa(51,32) -- \Aa(52,32);
        \draw[short] \Aa(52,32) -- \Aa(53,32);
        \draw[short] \Aa(53,32) -- \Aa(54,32);
        \draw[short] \Aa(54,32) -- \Aa(55,32);
        \draw[short] \Aa(55,32) -- \Aa(56,32);
        \draw[short] \Aa(56,32) -- \Aa(57,32);
        \draw[short] \Aa(57,32) -- \Aa(58,32);
        \draw[short] \Aa(58,32) -- \Aa(59,32);
        \draw[short] \Aa(59,32) -- \Aa(60,32);
        \draw[short] \Aa(60,32) -- \Aa(61,32);
        \draw[short] \Aa(61,32) -- \Aa(62,32);
        \draw[short] \Aa(62,32) -- \Aa(63,32);
        \draw[short] \Aa(63,32) -- \Aa(64,32);
        \draw[short] \Aa(64,32) -- \Aa(65,32);
        \draw[short] \Aa(65,32) -- \Aa(66,32);
        \draw[short] \Aa(66,32) -- \Aa(67,32);
        \draw[short] \Aa(67,32) -- \Aa(68,32);
        \draw[short] \Aa(68,32) -- \Aa(69,32);
        \draw[short] \Aa(69,32) -- \Aa(70,32);
        \draw[short] \Aa(70,32) -- \Aa(71,32);
        \draw[short] \Aa(71,32) -- \Aa(72,32);
        \draw[short] \Aa(72,32) -- \Aa(73,32);
        \draw[short] \Aa(73,32) -- \Aa(74,32);
        \draw[short] \Aa(74,32) -- \Aa(75,32);
        \draw[short] \Aa(75,32) -- \Aa(76,32);
        \draw[short] \Aa(76,32) -- \Aa(77,32);
        \draw[short] \Aa(77,32) -- \Aa(78,32);
        \draw[short] \Aa(78,32) -- \Aa(79,32);
        \draw[short] \Aa(79,32) -- \Aa(80,32);
        \draw[short] \Aa(80,32) -- \Aa(81,32);

        \def\Aa(#1,#2){(0.03125*#1-0.078125*#2+2,0.03125*#2+1)}
        
        \draw[short] \Aa(0,0) -- \Aa(0,1);
        \draw[short] \Aa(0,1) -- \Aa(0,2);
        \draw[short] \Aa(0,2) -- \Aa(1,2);
        \draw[short] \Aa(1,2) -- \Aa(1,3);
        \draw[short] \Aa(1,3) -- \Aa(1,4);
        \draw[short] \Aa(1,4) -- \Aa(2,4);
        \draw[short] \Aa(2,4) -- \Aa(3,4);
        \draw[short] \Aa(3,4) -- \Aa(4,4);
        \draw[short] \Aa(4,4) -- \Aa(4,5);
        \draw[short] \Aa(4,5) -- \Aa(4,6);
        \draw[short] \Aa(4,6) -- \Aa(5,6);
        \draw[short] \Aa(5,6) -- \Aa(5,7);
        \draw[short] \Aa(5,7) -- \Aa(5,8);
        \draw[short] \Aa(5,8) -- \Aa(6,8);
        \draw[short] \Aa(6,8) -- \Aa(7,8);
        \draw[short] \Aa(7,8) -- \Aa(8,8);
        \draw[short] \Aa(8,8) -- \Aa(9,8);
        \draw[short] \Aa(9,8) -- \Aa(10,8);
        \draw[short] \Aa(10,8) -- \Aa(11,8);
        \draw[short] \Aa(11,8) -- \Aa(12,8);
        \draw[short] \Aa(12,8) -- \Aa(12,9);
        \draw[short] \Aa(12,9) -- \Aa(12,10);
        \draw[short] \Aa(12,10) -- \Aa(13,10);
        \draw[short] \Aa(13,10) -- \Aa(13,11);
        \draw[short] \Aa(13,11) -- \Aa(13,12);
        \draw[short] \Aa(13,12) -- \Aa(14,12);
        \draw[short] \Aa(14,12) -- \Aa(15,12);
        \draw[short] \Aa(15,12) -- \Aa(16,12);
        \draw[short] \Aa(16,12) -- \Aa(16,13);
        \draw[short] \Aa(16,13) -- \Aa(16,14);
        \draw[short] \Aa(16,14) -- \Aa(17,14);
        \draw[short] \Aa(17,14) -- \Aa(17,15);
        \draw[short] \Aa(17,15) -- \Aa(17,16);
        \draw[short] \Aa(17,16) -- \Aa(18,16);
        \draw[short] \Aa(18,16) -- \Aa(19,16);
        \draw[short] \Aa(19,16) -- \Aa(20,16);
        \draw[short] \Aa(20,16) -- \Aa(21,16);
        \draw[short] \Aa(21,16) -- \Aa(22,16);
        \draw[short] \Aa(22,16) -- \Aa(23,16);
        \draw[short] \Aa(23,16) -- \Aa(24,16);
        \draw[short] \Aa(24,16) -- \Aa(25,16);
        \draw[short] \Aa(25,16) -- \Aa(26,16);
        \draw[short] \Aa(26,16) -- \Aa(27,16);
        \draw[short] \Aa(27,16) -- \Aa(28,16);
        \draw[short] \Aa(28,16) -- \Aa(29,16);
        \draw[short] \Aa(29,16) -- \Aa(30,16);
        \draw[short] \Aa(30,16) -- \Aa(31,16);
        \draw[short] \Aa(31,16) -- \Aa(32,16);
        \draw[short] \Aa(32,16) -- \Aa(32,17);
        \draw[short] \Aa(32,17) -- \Aa(32,18);
        \draw[short] \Aa(32,18) -- \Aa(33,18);
        \draw[short] \Aa(33,18) -- \Aa(33,19);
        \draw[short] \Aa(33,19) -- \Aa(33,20);
        \draw[short] \Aa(33,20) -- \Aa(34,20);
        \draw[short] \Aa(34,20) -- \Aa(35,20);
        \draw[short] \Aa(35,20) -- \Aa(36,20);
        \draw[short] \Aa(36,20) -- \Aa(36,21);
        \draw[short] \Aa(36,21) -- \Aa(36,22);
        \draw[short] \Aa(36,22) -- \Aa(37,22);
        \draw[short] \Aa(37,22) -- \Aa(37,23);
        \draw[short] \Aa(37,23) -- \Aa(37,24);
        \draw[short] \Aa(37,24) -- \Aa(38,24);
        \draw[short] \Aa(38,24) -- \Aa(39,24);
        \draw[short] \Aa(39,24) -- \Aa(40,24);
        \draw[short] \Aa(40,24) -- \Aa(41,24);
        \draw[short] \Aa(41,24) -- \Aa(42,24);
        \draw[short] \Aa(42,24) -- \Aa(43,24);
        \draw[short] \Aa(43,24) -- \Aa(44,24);
        \draw[short] \Aa(44,24) -- \Aa(44,25);
        \draw[short] \Aa(44,25) -- \Aa(44,26);
        \draw[short] \Aa(44,26) -- \Aa(45,26);
        \draw[short] \Aa(45,26) -- \Aa(45,27);
        \draw[short] \Aa(45,27) -- \Aa(45,28);
        \draw[short] \Aa(45,28) -- \Aa(46,28);
        \draw[short] \Aa(46,28) -- \Aa(47,28);
        \draw[short] \Aa(47,28) -- \Aa(48,28);
        \draw[short] \Aa(48,28) -- \Aa(48,29);
        \draw[short] \Aa(48,29) -- \Aa(48,30);
        \draw[short] \Aa(48,30) -- \Aa(49,30);
        \draw[short] \Aa(49,30) -- \Aa(49,31);
        \draw[short] \Aa(49,31) -- \Aa(49,32);
        \draw[short] \Aa(49,32) -- \Aa(50,32);
        \draw[short] \Aa(50,32) -- \Aa(51,32);
        \draw[short] \Aa(51,32) -- \Aa(52,32);
        \draw[short] \Aa(52,32) -- \Aa(53,32);
        \draw[short] \Aa(53,32) -- \Aa(54,32);
        \draw[short] \Aa(54,32) -- \Aa(55,32);
        \draw[short] \Aa(55,32) -- \Aa(56,32);
        \draw[short] \Aa(56,32) -- \Aa(57,32);
        \draw[short] \Aa(57,32) -- \Aa(58,32);
        \draw[short] \Aa(58,32) -- \Aa(59,32);
        \draw[short] \Aa(59,32) -- \Aa(60,32);
        \draw[short] \Aa(60,32) -- \Aa(61,32);
        \draw[short] \Aa(61,32) -- \Aa(62,32);
        \draw[short] \Aa(62,32) -- \Aa(63,32);
        \draw[short] \Aa(63,32) -- \Aa(64,32);
        \draw[short] \Aa(64,32) -- \Aa(65,32);
        \draw[short] \Aa(65,32) -- \Aa(66,32);
        \draw[short] \Aa(66,32) -- \Aa(67,32);
        \draw[short] \Aa(67,32) -- \Aa(68,32);
        \draw[short] \Aa(68,32) -- \Aa(69,32);
        \draw[short] \Aa(69,32) -- \Aa(70,32);
        \draw[short] \Aa(70,32) -- \Aa(71,32);
        \draw[short] \Aa(71,32) -- \Aa(72,32);
        \draw[short] \Aa(72,32) -- \Aa(73,32);
        \draw[short] \Aa(73,32) -- \Aa(74,32);
        \draw[short] \Aa(74,32) -- \Aa(75,32);
        \draw[short] \Aa(75,32) -- \Aa(76,32);
        \draw[short] \Aa(76,32) -- \Aa(77,32);
        \draw[short] \Aa(77,32) -- \Aa(78,32);
        \draw[short] \Aa(78,32) -- \Aa(79,32);
        \draw[short] \Aa(79,32) -- \Aa(80,32);
        


        \draw[short] (0,0) -- (2,0);

        \draw[color={rgb,255:red,100; green,200; blue,0}, dashed] (0,0) -- (0,2);

        \draw[color={rgb,255:red,100; green,200; blue,0}, dashed] (2,0) -- (2,2);

 
    \end{circuitikz}
    
    

        
    \caption{A Markov partition for $\begin{pmatrix}
        2 & 1 \\ 0 & 2
    \end{pmatrix}$.}
    \label{fig:hybrid MP}
\end{figure}

\begin{theorem}\label{thm:Jordans have hybrid MP}
    If $A$ is a $2\times 2$ expanding, non-diagonalizable integer matrix, then $f_A$ admits a hybrid Markov partition.
\end{theorem}

\begin{proof}
    Combine the construction from \cite{Be} with Lemma \ref{lem:MPforsimilarmatrices}.
\end{proof}

\section{Estimating Hausdorff dimension} \label{Section:FractalDim}


In this section, we estimate the Hausdorff dimension $\dim_H(\partial\M)$, using techniques from symbolic dynamics. We remark that some results of Urbanksi and Mañé \cite{U, Ma} can be applied to estimate the box dimension, or capacity, of $\partial\M$, and that for the explicit construction of Bedford \cite{Be}, some bounds from Dekking apply. Moreover, a recent result due to Ishii and Oka \cite{IO} gives a precise formula for $\dim_H(\partial\M)$, under a number of assumptions on the matrix inducing the endomorphism and the particular construction of $\M$. The bound given in Theorem \ref{thm-dimH} strongly resembles their formula when $A$ is conformal. The theorem is also in the spirit of the relation between Hausdorff dimension, entropy, and Lyapunov exponents developed by Manning and Young \cite{M1,Y}.


The argument here will be given in more generality. Recall from Section \ref{Section:Prelims} that for a positively expansive map $f:X\to X$ from a manifold onto itself, \cite{CR} shows we can find a compatible metric and  $\lambda>1,\varepsilon>0$ for which \begin{align*}d(x,y)<\varepsilon \implies d(fx,fy)\geq \lambda d(x,y).\tag{*}\end{align*} For this section, we will only work with a metric that is \textit{adapted} in this sense. In the case of expanding toral endomorphisms, this metric need not be the standard metric inherited from the Euclidean metric on $\R^d$. However, using the real Jordan form of an expanding integer matrix $A$, one can always find a metric for which (*) holds, that is also bi-Lipschitz equivalent to the standard metric. In particular, this means Hausdorff dimensions are preserved between the two metrics.  

Suppose we are given an expanding dynamical system $f:X\to X$ where $X$ is equipped with an adapted metric and $\M$ is a Markov partition for $f$. Any expanding dynamical system is positively expansive, with expansive constant $c=\varepsilon$. By refining sufficiently often, we may assume that for each rectangle $R\in \M$ we have that $\operatorname{diam}(R) < c/2$. Let, $T$ be the corresponding transition matrix and $\pi:\Sigma_T\to X$ be the SFT and covering map coming from $T$. Then, for any $x,y\in \Sigma_T$, we have that $$\pi(x)=\pi(y) \iff \forall i\in \N, R_{x_i}\cap R_{y_i}\neq \emptyset.$$

Mirroring a classical construction of Manning \cite{M2}, define a graph with vertex set $$\{(a,b): 1\leq a,b \leq |\M|, a\neq b, R_a\cap R_b \neq \emptyset\}$$ and an edge $(a,b)\to (c,d)$ if and only if $T_{ac}=T_{bd}=1$. Call the (one-sided) shift on this new graph $E_2$. We now define a map $\pi_2:E_2\to X$ by $$\pi_2((x_0,y_0),(x_1,y_1),...) = \pi(x_0x_1...).$$

It follows from Theorem 1.5 of \cite{AKS} that $\pi_2$ is surjective and for some $M$, $$\sup_{x\in \partial\M} |\pi_2^{-1}(x)|\leq M.$$ As a consequence, $\pi_2$ preserves topological entropy $h(\cdot)$, i.e. $$h(E_2)=h(f|_{\partial\M}).$$ For a purely topological proof of this fact, see Chapter IX, Theorem 1.8 of \cite{Ro}.


\begin{theorem}\label{thm-dimH}
    In the setup as above, $$\dim_H(\partial\M)\leq\frac{h(f|_{\partial\M})}{\log \lambda}.$$ 
\end{theorem}


\begin{proof} 

For a shift space $\Sigma$, write $B_n(\Sigma)$ for the \textit{allowed words} of length $n$ in $\Sigma$, that is, all sequences of length $n$ that appear in some point in $\Sigma$. Given a word $w$, let $C_w$ denote the cylinder set corresponding to the word in the appropriate subshift. Moreover, for a word $w$ allowed in $E_2$, let $w^1$ be the word formed by the first coordinates, and $w^2$ the word formed by the second coordinates.

First, for any $n\in \N$, 
$$\partial \M = \bigcup_{w\in B_n(E_2)} \pi_2(C_w) \subset \bigcup_{w\in B_n(E_2)} \pi(C_{w^1})\cup \pi(C_{w^2}).$$ We only need to consider words from $E_2$ since $f(\partial \M)\subset \partial \M$. Note that $C_w\subset E_2$, while $C_{w^1}$ and $C_{w^2}$ are subsets of $\Sigma_T$. 

By \cite{AKS} Lemma 3.2, there is some $D\in \R$ so that for any sufficiently long word $w'\in B_n(\Sigma_T)$, $$\operatorname{diam}(C_{w'}) \leq D \frac{1}{\lambda^n}.$$ 

Let $\mu\in \R$ be such that $h(E_2)=\log\mu$. Then $\# B_n(E_2) = E\mu^n + o(\mu^n)$ with $E\in \R$ by theorem 4.4.4 of \cite{LM}.


Hence, for fixed $s\in \R^+$, we have that \begin{align*}
\sum_{w\in B_n(E_2)} \operatorname{diam}(\pi_2(C_w))^s &\leq \sum_{w\in B_n(E_2)}\operatorname{diam}(\pi(C_{w^1})\cup \pi(C_{w^2}))^s \\
&\leq\sum_{w\in B_n(E_2)}\left(2D \frac{1}{\lambda^n}\right)^s  \\
&\leq \#B_n(E_2) \left(2D \frac{1}{\lambda^n}\right)^s \\
&= (2D)^s\frac{E\mu^n+o(\mu^n)}{\lambda^{ns}}
\end{align*}

As $n\to \infty$, this quantity goes to zero whenever $s>\frac{\log\mu}{\log\lambda}$, so $$\dim_H(\partial\M)\leq \frac{\log\mu}{\log\lambda} = \frac{h(E_2)}{\log\lambda}=\frac{h(f|_{\partial\M})}{\log\lambda},$$ as desired.

\end{proof}




It is possible to obtain equality: Consider the expanding toral endomorphism induced by $$A=\begin{pmatrix}
    2 & 0 \\ 0 & 2
\end{pmatrix}$$ on $\T^2$, with the Markov partition given by cubes of side length $1/5$. In this case, the adapted metric can be taken to be the standard metric, and $h(E_2)$ can be computed manually. Generally, the bound gets worse when the dimension or the difference in modulus between the eigenvalues of $A$ increases. 


Staying with expanding toral endomorphisms in dimension 2, our result implies that the boundary of any Markov partition must be non-trivial in the sense of entropy.

\begin{corollary}
    For any Markov partition for an expanding toral endomorphism on $\T^2$, $h(f|_{\partial\M})>0$.
\end{corollary}

\begin{proof}
    Since $\dim_H(\partial\M)\geq \dim_{top}(\partial\M) = 1$, Theorem \ref{thm-dimH} gives $$h(f|_{\partial\M})\geq \log\lambda > 0.$$
\end{proof}


\begin{remark}
    With some small adjustments, the approach to computing Hausdorff dimension in Theorem \ref{thm-dimH} can be extended to the invertible setting. The definition of a Markov partition is more complex, and this setting will be reserved for future work.
\end{remark}








\newpage

\begin{thebibliography}{WWWWW}

    \bibitem[AKS]{AKS} J. Ashley, B. Kitchens, and M. Stafford, {\em Boundaries of Markov partitions}, Trans. Amer. Math. Soc. {\bf 333}, 1 (1992), 177-201.

    \bibitem[AW]{AW} R. L. Adler and B. Weiss, {\em Similarity of automorphisms of the torus}, Mem. Amer. Math, Soc.,(1970).

    \bibitem[Be]{Be} T. Bedford, {\em Generating special Markov partitions using fractals}, Ergod. Th.  Dynam. Sys. {\bf 6} (1986), 325-333. 
    
    \bibitem[Bo1]{Bo1} R. Bowen, {\em Equilibrium States and the Ergodic Theory of Anosov Diffeomorphisms}, 2nd ed., Lect. Notes Math. {\bf 470} (2008). 

    \bibitem[Bo2]{Bo2} R. Bowen, {\em Markov partitions are not smooth}, Proc. Amer. Math. Soc. {\bf 71} (1978), 130-132.

    \bibitem[Bo3]{Bo3} R. Bowen, {\em Markov Partitions and Minimal Sets for Axiom A Diffeomorphisms}, Amer. Jour. Math. (1970), 907-918.
    
    \bibitem[C]{C} E. Cawley, {\em Smooth Markov partitions and toral automorphisms}, Ergod. Th. Dynam. Syst. {\bf 11} (1991), 633-651.

    \bibitem[CR]{CR} E. Coven and W. Reddy, {\em Positively expansive maps of compact manifolds}.  In: Global Theory of Dynamical Systems, Ed: Nitecki, Z., Robinson, C.. Lect. Notes in Math. {\bf 819} (1980). 

    \bibitem[D]{D} F. Dekking, {\em Recurrent sets}, Advances in Math. {\bf 44} (1982), 78-104.

    \bibitem[F]{F} J. M. Franks, {\em Invariant sets of hyperbolic toral automorphisms}, Amer. J. Math. {\bf 99}, 5 (1977), 1089-1095.
    
    \bibitem[IO]{IO} Y. Ishii and T. Oka, {\em On the Hausdorff dimension of the recurrent sets induced from endomorphisms of free groups},
    J. Fractal Geom. {\bf 9} (2022), 171–192. 

    \bibitem[K]{K} R. Koo, {\em A Classification of Matrices of Finite Order over $\C$, $\R$, and $\Q$}, Math. Mag. {\bf 76}, 2 (2003), 143-148.

    \bibitem[KP]{KP} J. Kuzmanovich and A. Pavlichenkov, {\em Finite Groups of Matrices Whose Entries Are Integers}, Amer. Math. Month. {\bf 109}, 2 (2002), 173-186.

    \bibitem[LM]{LM} D. Lind and B. Marcus, {\em Introduction to Symbolic Dynamics and Coding}, 2nd ed. (2021).

    \bibitem[Ma]{Ma} R. Mañé, {\em Orbits of paths under hyperbolic toral automorphisms}, Proc. Amer. Math. Soc. {\bf 73}, 1 (1979), 121-125.

    \bibitem[M1]{M1} A. Manning, {\em A relation between Lyapunov exponents, Hausdorff dimension and entropy}, Ergod. Th. Dynam. Sys. {\bf 1}, 4 (1981), 451-459.

    \bibitem[M2]{M2} A. Manning, {\em Axiom A diffeomorphisms have rational zeta functions}, Bull. Lon. Math. Soc. {\bf 3}, 2 (1971), 215-220.

    \bibitem[Ro]{Ro} C. Robinson. {\em Dynamical Systems: Stability, Symbolic Dynamics and Chaos}, (1995).
   
    \bibitem[Ru]{Ru} D. Ruelle. {\em Thermodynamic Formalism}, 2nd ed. (2004).
    
    \bibitem[Si]{Si} Ya. G. Sinai, {\em Construction of markov partitions},  Func. Ana. App. {\bf 2}, 3 (1968), 245-253.

    \bibitem[U]{U} M. Urbánski, {\em On the capacity of a continuum with non-dense orbit under a hyperbolic toral automorphism}, Studia Math. {\bf 81} (1985), 37-51.

    \bibitem[URM]{URM} M. Urbánski, M. Roy, and S. Munday, {\em Non-Invertible Dynamical Systems. Volume 1: Ergodic Theory – Finite and Infinite, Thermodynamic Formalism, Symbolic Dynamics and Distance Expanding Maps}, (2022). 

    \bibitem[Y]{Y} L.-S. Young, {\em Dimension, entropy and Lyapunov exponents}, Ergod. Th. Dynam. Sys. {\bf 2}, 1 (1982), 109-124.


	
\end{thebibliography}	


\end{document}
